%Version 3.1 December 2024
% See section 11 of the User Manual for version history
%
%%%%%%%%%%%%%%%%%%%%%%%%%%%%%%%%%%%%%%%%%%%%%%%%%%%%%%%%%%%%%%%%%%%%%%
%%                                                                 %%
%% Please do not use \input{...} to include other tex files.       %%
%% Submit your LaTeX manuscript as one .tex document.              %%
%%                                                                 %%
%% All additional figures and files should be attached             %%
%% separately and not embedded in the \TeX\ document itself.       %%
%%                                                                 %%
%%%%%%%%%%%%%%%%%%%%%%%%%%%%%%%%%%%%%%%%%%%%%%%%%%%%%%%%%%%%%%%%%%%%%

%%\documentclass[referee,sn-basic]{sn-jnl}% referee option is meant for double line spacing

%%=======================================================%%
%% to print line numbers in the margin use lineno option %%
%%=======================================================%%

%%\documentclass[lineno,pdflatex,sn-basic]{sn-jnl}% Basic Springer Nature Reference Style/Chemistry Reference Style

%%=========================================================================================%%
%% the documentclass is set to pdflatex as default. You can delete it if not appropriate.  %%
%%=========================================================================================%%

%%\documentclass[sn-basic]{sn-jnl}% Basic Springer Nature Reference Style/Chemistry Reference Style

%%Note: the following reference styles support Namedate and Numbered referencing. By default the style follows the most common style. To switch between the options you can add or remove “Numbered” in the optional parenthesis. 
%%The option is available for: sn-basic.bst, sn-chicago.bst%  
 
%%\documentclass[pdflatex,sn-nature]{sn-jnl}% Style for submissions to Nature Portfolio journals
%%\documentclass[pdflatex,sn-basic]{sn-jnl}% Basic Springer Nature Reference Style/Chemistry Reference Style
\documentclass[pdflatex,sn-mathphys-num]{sn-jnl}% Math and Physical Sciences Numbered Reference Style
%%\documentclass[pdflatex,sn-mathphys-ay]{sn-jnl}% Math and Physical Sciences Author Year Reference Style
%%\documentclass[pdflatex,sn-aps]{sn-jnl}% American Physical Society (APS) Reference Style
%%\documentclass[pdflatex,sn-vancouver-num]{sn-jnl}% Vancouver Numbered Reference Style
%%\documentclass[pdflatex,sn-vancouver-ay]{sn-jnl}% Vancouver Author Year Reference Style
%%\documentclass[pdflatex,sn-apa]{sn-jnl}% APA Reference Style
%%\documentclass[pdflatex,sn-chicago]{sn-jnl}% Chicago-based Humanities Reference Style

%%%% Standard Packages
%%<additional latex packages if required can be included here>

\usepackage[utf8]{inputenc}
\usepackage{graphicx}%
\usepackage{multirow}%
\usepackage{amsmath,amssymb,amsfonts}%
\usepackage{amsthm}%
\usepackage{mathrsfs}%
\usepackage[title]{appendix}%
\usepackage{xcolor}%
\usepackage{textcomp}%
\usepackage{manyfoot}%
\usepackage{bookmark}
\usepackage{booktabs}%
\usepackage{array}
\usepackage{float}
\usepackage{placeins}
\usepackage[utf8]{inputenc}
\usepackage{newunicodechar}
\newunicodechar{（}{(}
\newunicodechar{）}{)}
\usepackage{array}
\usepackage{svg}
\usepackage{caption}
\usepackage{subcaption}
\usepackage{algorithm}%
\usepackage{algorithmicx}%
\usepackage{algpseudocode}%
\usepackage{listings}%

\usepackage{dblfloatfix}
\usepackage{url}
%%%%

%%%%%=============================================================================%%%%
%%%%  Remarks: This template is provided to aid authors with the preparation
%%%%  of original research articles intended for submission to journals published 
%%%%  by Springer Nature. The guidance has been prepared in partnership with 
%%%%  production teams to conform to Springer Nature technical requirements. 
%%%%  Editorial and presentation requirements differ among journal portfolios and 
%%%%  research disciplines. You may find sections in this template are irrelevant 
%%%%  to your work and are empowered to omit any such section if allowed by the 
%%%%  journal you intend to submit to. The submission guidelines and policies 
%%%%  of the journal take precedence. A detailed User Manual is available in the 
%%%%  template package for technical guidance.
%%%%%=============================================================================%%%%

%% as per the requirement new theorem styles can be included as shown below
\theoremstyle{thmstyleone}%
\newtheorem{theorem}{Theorem}%  meant for continuous numbers
%%\newtheorem{theorem}{Theorem}[section]% meant for sectionwise numbers
%% optional argument [theorem] produces theorem numbering sequence instead of independent numbers for Proposition
\newtheorem{proposition}[theorem]{Proposition}% 
%%\newtheorem{proposition}{Proposition}% to get separate numbers for theorem and proposition etc.

\theoremstyle{thmstyletwo}%
\newtheorem{example}{Example}%
\newtheorem{remark}{Remark}%

\theoremstyle{thmstylethree}%
\newtheorem{definition}{Definition}%

\raggedbottom
%%\unnumbered% uncomment this for unnumbered level heads

\begin{document}

\title[Article Title] {\textcolor{blue}{MGAT-FinRL: A Multi-relational Graph Attention Transformer-Based Reinforcement Learning Investment Strategy}}

%%=============================================================%%
%% GivenName	-> \fnm{Joergen W.}
%% Particle	-> \spfx{van der} -> surname prefix
%% FamilyName	-> \sur{Ploeg}
%% Suffix	-> \sfx{IV}
%% \author*[1,2]{\fnm{Joergen W.} \spfx{van der} \sur{Ploeg} 
%%  \sfx{IV}}\email{iauthor@gmail.com}
%%=============================================================%%

\author*[1]{\fnm{Zhiping} \sur{Chen}}\email{zchen@mail.xjtu.edu.cn}

\author[1]{\fnm{Xiaoping} \sur{Ren}}\email{xpren0703@stu.xjtu.edu.cn}

\affil*[1]{\orgdiv{School of Mathematics and Statistics}, \orgname{Xi’an Jiaotong University}, \orgaddress{\city{Xi’an}, \postcode{710049}, \state{Shaanxi}, \country{China}}}

%%==================================%%
%% Sample for unstructured abstract %%
%%==================================%%

\abstract{Quantitative investment aims to identify an optimal security trading strategy, which highlights the critical role of market information in assessing price trends. However, most existing strategies rely primarily on asset-level trading data and neglect cross-firm dependencies arising from shared ownership, investor-relations disclosures, and correlated trading activity. {\color{blue}To address this gap, we propose MGAT-FinRL, which combines three relation graphs (common top-ten shareholders, co-mentions in investor-relations questions, and Kendall correlations of trading value) with a multi-relational graph-attention Transformer (MGAT) that forecasts the next two days' high, low, and closing prices and is trained to rank the return available under next-day execution. Its out-of-sample forecasts enter the state of five deep reinforcement learning agents (PPO, A2C, SAC, TD3, and DDPG) at every decision. On fifteen-stock universes from the Chinese banking and healthcare sectors, the agents were trained on 2022--2023, selected on 2024, and evaluated once on 2025 under a protocol registered before the 2025 data were collected, with next-day execution, trading costs, five paired seeds, and block-bootstrap inference. In 2025, MGAT forecasts the next day's high and low prices 15--29\% more accurately than previous-day values, and its forecasts of the tradable-horizon return have a significant rank correlation with realised returns in healthcare (0.052, $p=0.012$), but not in banking (0.014). Agents that observe the MGAT forecasts do not attain a significantly higher Sharpe ratio than agents without forecast information (paired differences of 0.04 in healthcare and $-0.04$ in banking), and their allocations remain close to equal weighting and nearly insensitive to the forecasts. A ranking rule extracts the forecast information before costs (Sharpe ratio 1.48 versus 1.25 for equal weighting in healthcare) but not after. These results locate the main constraint of the forecast-to-policy pipeline at the interface between a weak cross-sectional signal and the policy learner.}}

\keywords{\textcolor{blue}{graph attention network, Transformer, deep reinforcement learning, portfolio allocation, financial relation graphs, pre-registered evaluation}}

%%\pacs[JEL Classification]{D8, H51}

%%\pacs[MSC Classification]{35A01, 65L10, 65L12, 65L20, 65L70}

\maketitle

\section{Introduction}\label{sec1}

Financial markets are increasingly shaped by rapid information diffusion, cross-asset spillovers, and abrupt regime shifts, making robust quantitative trading systems more important and more difficult than ever. Quantitative investment aims to translate heterogeneous signals into disciplined decisions, mitigating behavioral bias while improving risk-adjusted performance. In practice, however, the current prediction-to-decision pipeline faces two persistent obstacles: (i) market observations are noisy and non-stationary, and (ii) price dynamics are not only temporal but also relational, driven by evolving connections among firms, investors, and industries.

Prior research in this area has broadly progressed along two complementary directions. The first direction focused on improving forecasting models to suppress noise and enhance prediction accuracy, including graph-based learners and deep sequence architectures. The second direction emphasized investment strategy design and portfolio optimization, often via reinforcement learning (RL) and other decision-making paradigms. Although both directions have advanced substantially, they were often developed in relative isolation, which limited practical deployment: a superior forecaster does not necessarily yield a strong trading strategy, and a powerful decision policy can be undermined by weak or unstable predictive signals. This motivates an integrated view that jointly considers representation learning, forecasting, and portfolio construction in a unified framework. In this regard, three closely related challenges are central, as shown below.

\paragraph{Relational feature extraction}
 Stock prices are influenced by a wide range of factors, spanning macroeconomic conditions, sector dynamics, corporate fundamentals, trading behaviours, news, and governance. 
Graph neural networks (GNNs) have provided effective tools for capturing such latent dependencies, including temporal graph forecasting~\citep{Zhao2020TGCN}, 
spatio-temporal formulations~\citep{Jafari2022GCNET,Feng2025STGAT}, 
and heterogeneous relational modelling~\citep{Zhao2024RAHGNN,Sun2025GTC,Wu2024MomentumKG}, 
with adaptive graph filtering further improving robustness~\citep{Andrea2023AdaptiveFiltersGCNN}. 
However, most existing graph constructions primarily rely on quantitative co-movements or coarse metadata, and they rarely operationalize qualitative relations such as shareholder linkages and investor-relations interactions in a systematic, learnable manner\textcolor{blue}{, which motivates richer multi-source relation modelling}.
{\color{blue}Combining complementary representations is a recurring design principle in deep learning. In medical image analysis, for example, multi-scale region-based convolutional networks~\citep{Bakkouri2019MSCNN}, multi-layer feature fusion networks~\citep{Bakkouri2020CAD}, bidirectional gated 3D multi-scale fusion~\citep{Bakkouri2022BG3DM2F}, and multi-level context attentional fusion~\citep{Bakkouri2023MLCA2F} combine features from different receptive fields or network depths. MGAT applies a related principle to financial relations by combining shareholding, investor-relations Q\&A, and trading-activity graphs within a single graph-attention encoder.\par}

\paragraph{Price prediction with relational structure}
Once relational signals are represented, the next challenge is to forecast prices effectively under non-stationary dynamics. 
A considerable body of work has proposed deep-learning architectures for stock price trend and closing-price prediction, including multi-period CNN-GRU frameworks and attention-based recurrent models
\citep{Luo2023MethanolFutures,Lee2022GRUAttentionTrading,Zhang2022TSLALSTMGRU},
as well as graph-enhanced predictors and hybrid graph-sequence designs that explicitly model inter-stock relations
\citep{Ye2020MultiGraph,Yun2019GraphTransformerNetworks,Li2022FDGRUTransformer,Lu2025BiLSTMGAT}.
These approaches often performed well for short-horizon movements, yet \textcolor{blue}{forecasting from multi-day histories} requires models that reconcile (i) local relational inductive bias and (ii) global temporal context over extended horizons. 
To this end, transformer-style self-attention has been widely adopted to capture long-range temporal dependencies in noisy sequences, while graph transformer variants further extend attention to relational settings by improving global interaction modelling beyond locality-constrained message passing
\citep{Shehzad2024GraphTransformersSurvey,Gao2024DiffsFormer,Wu2023NodeFormer,Shirzad2023Exphormer,Chen2022StructureAwareTransformerGraphs}.

However, many advances were primarily examined from a representation-learning perspective and did not explicitly target \textcolor{blue}{predictive performance under real-world trading constraints; composite architectures that couple GNN-based relational encoding with temporal self-attention along the input horizon have rarely been evaluated in this setting}.

\paragraph{Efficient portfolio construction}
Beyond prediction, portfolio construction is a sequential decision-making problem that must manage risk, transaction costs, and changing market conditions. 
RL-based portfolio selection methods have demonstrated strong potential, including multi-module deep RL for trading \citep{Ma2021ParallelMultiModuleDRLTrading}, the FinRL framework for reproducible financial RL pipelines \citep{Liu2021FinRL}, deep graph-convolutional RL for relationally informed decision making \citep{Farzan2021DeepPocket}, and deep differentiable RL formulations that enable end-to-end optimisation through differentiable components \citep{Thibault2022DeepDifferentiableRLTrading}.
However, existing approaches are often trained predominantly on market-derived signals, while RL evaluations frequently vary in state design, investment constraints, and training protocols, making the incremental value of forecast-driven signals difficult to quantify under a controlled setting. 

To address these issues, we propose a multi-relational graph attention Transformer-based reinforcement learning investment strategy (\emph{MGAT-FinRL}), a forecast-to-policy framework with a standardised interface between the forecaster and the decision layer and an evaluation protocol fixed before the evaluation data were collected.
The main contributions are summarised as follows:
\begin{itemize}

  \item \textbf{Multi-layer relational graph construction from heterogeneous financial ties.}
  We model stock dependencies with three complementary relation layers, including (i) an investor-interaction graph derived from the board-secretary investor-relations Q\&A disclosures, (ii) an equity-holder relation graph capturing shareholder linkages, and (iii) a daily \textcolor{blue}{trading-value} co-movement graph based on Kendall \textcolor{blue}{rank} correlations, forming a structured relational prior for downstream learning.

  \item \textbf{\textcolor{blue}{A graph-attention Transformer forecaster trained for the tradable horizon.}}
  {\color{blue}MGAT applies graph attention restricted to the edges of the three relation graphs to every day of a 20-day window and models each stock's sequence with a Transformer encoder. It is trained with a squared-error loss plus a ranking term for the return from the close of $t+1$ to the close of $t+2$, the first return available to a decision made at the close of $t$. We evaluate it in four walk-forward folds (2022--2025) against previous-day values, a ridge regression, a graph-free Transformer, and three graph encoders (graph convolution, no cross-stock propagation, and the multi-relation encoder of the original submission).\par}

  \item {\color{blue}\textbf{A pre-registered forecast-to-policy evaluation on an unseen year.}
  Five DRL agents receive the forecasts in their state at every decision and are trained on 2022--2023, selected on 2024, and evaluated once on 2025, a year that entered no earlier analysis; the hypotheses, sources, agents, seeds, and tests were committed to a public repository before the 2025 data were downloaded. Comparing agents that differ only in the forecast source (MGAT, graph-free, ridge, and a naive source without forecasts) under next-day execution, trading costs, and paired block-bootstrap inference measures the contribution of the forecaster to the learned allocations. In 2025, the MGAT forecasts rank the tradable return significantly in healthcare, whereas the agents that observe them do not attain a significantly higher Sharpe ratio than agents without forecasts; policy diagnostics, cost sensitivity, an encoder ablation, and forecast-ranking rules show where the transfer from forecasts to allocations is constrained.\par}
\end{itemize}

\section{Related work}\label{sec2}

We have seen numerous AI-based studies on security price forecasts and quantitative investment strategies. Among these research areas, GNNs have demonstrated their strong capabilities in capturing dynamic relationships among financial instruments by modelling the market as a graph, thereby enabling context-aware feature representations. Concurrently, transformer-based models have achieved state-of-the-art performance in time-series forecasting, owing to their self-attention mechanism and capacity to capture long-term dependencies in sequential financial data. Last but not least, reinforcement learning frameworks such as FinRL \citep{Liu2021FinRL} offer a flexible and data-driven approach for training trading agents that can learn optimal portfolios through continuous interaction with the market environment. \textcolor{blue}{Combining variants of these techniques makes it possible to address the three challenges outlined above within one pipeline.} In view of this, we will mainly review related works about these three AI techniques in what follows.

\paragraph{GNNs for relationship stock modelling}
GNNs were introduced to learn representations from graph-structured data by propagating and aggregating information along edges \citep{Scarselli2009GNNModel}. 
Convolutional GNN variants (ConvGNNs) made neighbourhood aggregation more scalable, which facilitated their adoption in quantitative finance for modelling inter-stock relations \citep{Chen2018CIKMStockGCN}.
Early and representative efforts typically cast the market as correlation- or relation-driven graphs and learn context-aware node features for prediction. 
Graph-Based Prediction Using Graph Convolutional Network (GCNET) captured inter-stock dynamics by explicitly modelling multiple relations among stocks \citep{Jafari2022GCNET}, while auxiliary or heterogeneous information was incorporated to enrich relational signals beyond pure price co-movement \citep{Andrea2023AdaptiveFiltersGCNN}. 
More recently, heterogeneous and attention-based designs such as GNN-transformer co-contrastive learning (GTC) provided a unified mechanism to contextualize node representations on heterogeneous graphs \citep{Sun2025GTC}. 
Notably, a multi-graph convolution adversarial network (VGC-GAN) built multiple correlation graphs from historical data and combined multi-graph convolution with adversarial training to enhance prediction robustness \citep{Ma2024VGC_GAN}.
However, many financial GNN pipelines still rely primarily on trading-derived relations or coarse metadata, while rich qualitative relations remain comparatively less explored in a unified graph-learning pipeline.

\paragraph{Transformers and graph-aware transformers for forecasting}
Beyond relational modelling, \textcolor{blue}{forecasting} critically depends on capturing temporal dependencies in noisy and non-stationary price sequences. 
Transformers have achieved strong performance in time-series forecasting due to their self-attention and the ability to model long-range temporal interactions. 
In finance, quite a few studies empirically showed that transformer-based modelling can improve index prediction and backtesting performance across major markets \citep{Wang2022ESWA_TransformerIndex}. 
However, most of these successes were achieved in a pure sequence setting, whereas relational forecasting in markets additionally requires encoding graph structure, propagating information beyond local neighbourhoods, and remaining computationally efficient.
For this reason, graph-orientated transformer variants explored these issues from complementary perspectives. 
Ying et al. introduced structural encodings for transformers on graphs \citep{Ying2021NeurIPSTransformersGraph}, while Graph Transformer Networks (GTNs) composed meta-path-based relations to capture heterogeneous dependencies \citep{Yun2019GraphTransformerNetworks}. 
Subsequent work improved expressiveness and scalability via techniques such as diffusion-style propagation \citep{Gao2024DiffsFormer}, all-pairs message passing \citep{Wu2023NodeFormer}, sparse attention designs \citep{Shirzad2023Exphormer}, and structure-aware encoding \citep{Chen2022StructureAwareTransformerGraphs}. 
Surveys further systematised the design space and empirical findings of Graph Transformer \citep{Shehzad2024GraphTransformersSurvey}.
However, existing transformer and graph-transformer studies in finance mainly emphasise representation learning or short-term improvements, without clearly operationalising long-term temporal dependency modelling in a way that is compatible with multi-asset trading scenarios and controlled decision evaluation.
{\color{blue}Related advances in time-series classification have strengthened temporal representation learning through knowledge distillation and capsule-based modelling. Xiao et al.~\citep{Xiao2026ETBDSD} proposed an ensemble transitive bidirectional decoupled self-distillation scheme that transfers averaged multi-level features between higher and lower network levels, and Xiao et al.~\citep{KASSCN2026} proposed KACapMamba, a knowledge-aware state-space capsule network that combines attentive Mamba blocks with capsule routing for multivariate time-series classification. These methods classify labelled sequences, whereas MGAT learns cross-stock relational representations for price forecasting and supplies its forecasts to a downstream portfolio policy.\par}

\paragraph{RL for portfolio selection}
While GNNs and transformers mainly improve predictive representations by exploiting relational structure and long-range temporal dependencies, a deployable quantitative system ultimately requires a decision layer that converts such signals into executable investment actions under constraints. 
Portfolio selection is inherently sequential and must account for risk, transaction costs, and changing market conditions; reinforcement learning (RL) has therefore been increasingly adopted for trading and portfolio management. 
Representative works included multi-module deep RL designs \citep{Ma2021ParallelMultiModuleDRLTrading} and open frameworks such as FinRL that enabled data-driven training and benchmarking of trading agents \citep{Liu2021FinRL}. 
Other studies proposed richer architectures to improve representation and policy learning, such as deep graph convolutional RL for portfolio management \citep{Farzan2021DeepPocket} and differentiable RL formulations for optimal trading \citep{Thibault2022DeepDifferentiableRLTrading}. 
More recent directions emphasized transparency and robustness, e.g., explainable deep RL for portfolio management \citep{Guan2021ExplainableDRLPortfolio} and model-free deep Q-learning variants for execution \citep{Brian2021DoubleDeepQLearning}.
However, many RL pipelines were trained primarily on market time-series data with limited integration of heterogeneous relational and qualitative context, and RL evaluations also vary substantially in state design and constraint modelling.

Overall, GNNs improved relational feature extraction for stock prediction, transformers strengthened long-term temporal modelling, and RL provided adaptive investment decision mechanism. 
However, these components are still rarely integrated under a controlled and reproducible forecast-to-policy pipeline; therefore, a unified framework that jointly models heterogeneous relations, performs \textcolor{blue}{forecasting with explicit temporal dependence over a multi-day input window}, and translates forecasts into portfolio selection \textcolor{blue}{remains to be evaluated under a controlled protocol}. \textcolor{blue}{We build a pipeline of this kind in what follows and evaluate it under a controlled protocol.}

{\color{blue}\paragraph{Positioning of this study}
Closest to our work are graph-based DRL portfolio models such as GPM~\citep{Si2022GPM}, multi-GCGRU~\citep{Liu2021MultiGCGRU}, and DeepPocket~\citep{Farzan2021DeepPocket}, which embed relational information about stocks into the portfolio decision model, and graph-attention forecasters such as HATS~\citep{Kim2019HATS}, STGAT~\citep{Feng2025STGAT}, and BiLSTM-GAT~\citep{Lu2025BiLSTMGAT}, which weight related stocks by learned attention. MGAT-FinRL differs in three respects. First, its relation layers are built from qualitative financial ties (common top-ten shareholders and co-mentions in investor-relations questions) in addition to a trading-activity correlation, and they restrict, rather than weight, the attention between stocks; in the present data, the trading-value layer determines most of the admissible pairs (Section~\ref{subsec3}). Second, the forecaster is trained for the return that a decision can actually earn under next-day execution, and its forecasts enter the policy's state at every decision; forecasting and policy learning remain separate stages, so the forecast source can be varied while the policy optimiser, environment, and seeds are held fixed. Third, the evaluation is pre-registered and uses a year that entered no earlier analysis, with a naive forecast source of identical state dimension as the reference. The individual GAT, Transformer, and actor-critic components are standard; the contribution lies in the relation construction, the tradable-horizon objective, and a controlled evaluation of how much relational forecasts change portfolio outcomes.\par}

\section{Architecture of MGAT-FinRL}\label{sec3}

{\color{blue}MGAT-FinRL links three components that are trained in sequence: a set of relation graphs built from shareholding, investor-relations, and trading-activity data; a forecaster that propagates daily stock features over these graphs with attention and then models each stock's recent history with a Transformer encoder; and a portfolio policy learned with deep reinforcement learning (DRL) whose observations contain the forecasts at every decision (Fig.~\ref{fig:fig1}). Multi-relation graph models such as the multi-GCN of \citep{Ye2020MultiGraph}, GPM \citep{Si2022GPM}, and multi-GCGRU \citep{Liu2021MultiGCGRU} weight the neighbours of a stock with fixed or globally learned coefficients. Graph attention \citep{Velickovic2018GAT} instead weights each neighbour according to the current features of both stocks, so the influence of a related firm can change from day to day while the set of admissible neighbours remains fixed by the relation graphs.\par}

\paragraph{Step 1: Relation graphs}
{\color{blue}Three static relation matrices are constructed for each sector: common top-ten shareholders ($\mathbf{A}_{\mathrm{SH}}$), co-mentions of two firms in the same investor-relations question ($\mathbf{A}_{\mathrm{QA}}$), and Kendall rank correlations of daily trading values ($\mathbf{A}_{\mathrm{KD}}$) (Section~\ref{Heterogeneous relation graphs}). Their positive entries are scaled to a common range and combined into one edge set, which defines the stocks that each stock may attend to.\par}

\paragraph{Step 2: Graph-attention Transformer forecaster (MGAT)}
{\color{blue}At the close of day $t$, MGAT receives a window of the 20 most recent trading days of eight features per stock (daily return, high and low relative to the close, relative trading value, and four valuation yields). Three graph-attention layers are applied to every day of the window, the resulting representation is concatenated with the inputs, and one Transformer encoder block per stock models the 20-day sequence. A linear layer maps the last-day representation to the high, low, and closing prices of days $t+1$ and $t+2$, expressed as log ratios to the close of $t$. Training combines the squared forecast error with a ranking term for the return from the close of $t+1$ to the close of $t+2$, the first return that a decision made at the close of $t$ can earn under next-day execution (Section~\ref{subsec:objective}).\par}

\paragraph{Step 3: Policy learning with forecasts in the state}
{\color{blue}The fitted forecaster is held fixed, and its out-of-sample forecasts are appended to the observation of each DRL agent at every decision: the cross-sectional $z$-score of the predicted tradable-horizon return and the standardised predicted intraday range of the next day. The remaining state components are those of the FinRL portfolio environment \citep{Liu2021FinRL}: a 252-day return covariance matrix and four technical indicators. Five agents (PPO, A2C, SAC, TD3, and DDPG) map the state to long-only portfolio weights, and all rewards and portfolio values are computed from observed prices net of trading costs (Section~\ref{sec4}). The two-stage design keeps the forecaster identical across agents and allows the same environment to be run with forecasts from different sources, including a naive source without forecast information, so that the contribution of the forecaster to allocation can be measured directly.\par}

In Fig.~\ref{fig:fig1}, $T$ denotes the input window length, $N$ the number of stocks (nodes), and $F$ the number of input features; the input is $\mathbf{X}_t\in\mathbb{R}^{T\times N\times F}$, and the graph-attention output is $\mathbf{H}^{\mathrm{GAT}}_t\in\mathbb{R}^{T\times N\times C}$.
\begin{sidewaysfigure}
  \centering
  \includegraphics[width=\textheight]{fig1_mgat.pdf}
  \caption{{\color{blue}Architecture of MGAT-FinRL. Upper part: the MGAT forecaster, in which graph attention restricted to the union of the three relation graphs is applied to every day of the input window and a Transformer encoder models each stock's 20-day sequence; the model is trained with a squared-error loss and a ranking term for the tradable-horizon return. Lower part: the FinRL module, in which the out-of-sample forecasts enter the state of each DRL agent at every decision.}}
  \label{fig:fig1}
\end{sidewaysfigure}

\section{MGAT forecaster}
\label{title_page}
Prior work has shown that graph-based modelling can improve financial prediction by explicitly encoding inter-asset dependencies and that relational priors can be integrated into DRL for portfolio management.
\citep{Ye2020MultiGraph,Si2022GPM,Liu2021MultiGCGRU,Sun2025GTC}.
However, many existing systems rely on a single relation type (e.g., industry or correlation graph) or are trained predominantly on trading-derived signals, leaving heterogeneous qualitative relations under-exploited and weakening robustness under event-driven or governance-driven regimes.
These limitations motivate the design of the MGAT forecaster within MGAT-FinRL, which couples \textcolor{blue}{graph attention over multiple relation graphs with a Transformer encoder for price forecasting}.

\subsection{Problem definition}
As shown in the upper part of Figure~\ref{fig:fig1}, \textcolor{blue}{MGAT} integrates \textcolor{blue}{graph-attention} relational encoding of listed companies and their trading attributes and further employs transformer self-attention \textcolor{blue}{to forecast the prices of each stock}.
To highlight the generality and improvements of \textcolor{blue}{MGAT} over existing architectures, we first recall the formulation of a representative multi-GCN model in \citep{Ye2020MultiGraph}:
\begin{equation}
\widehat{\mathbf{Y}}_{D}
=
f\!\left(
\left[
\mathbf{x}_{D-t},
\mathbf{x}_{D-t+1},
\ldots,
\mathbf{x}_{D-1},
\mathbf{v}_{D-t},
\mathbf{v}_{D-t+1},
\ldots,
\mathbf{v}_{D-1}
\right],
\mathcal{I},
\Theta
\right),
\label{eq:multigcn}
\end{equation}
where $\widehat{\mathbf{Y}}_{D}\in\mathbb{R}^{N\times 4}$
denotes the predicted OHLC price matrix for the $N$ stocks on
trading day $D$. Specifically, its $i$-th row is defined as
\[
\widehat{\mathbf{Y}}_{D}(i,:)
=
\left[
\widehat{p}^{\,\mathrm{open}}_{D,i},
\widehat{p}^{\,\mathrm{high}}_{D,i},
\widehat{p}^{\,\mathrm{low}}_{D,i},
\widehat{p}^{\,\mathrm{close}}_{D,i}
\right],
\]
where the four columns correspond to the predicted opening, highest,
lowest, and closing prices, respectively. Here, $D$ denotes the target
trading day, $\mathbf{x}_{D-k}$ represents the observed trading-feature
matrix $k$ days before day $D$, and $\mathbf{v}_{D-k}$ denotes the
corresponding valuation and fundamental-feature matrix, including
indicators such as PE, PB, PS, and PCF. In addition, $\mathcal{I}$
represents external relational information, and $\Theta$ denotes the
trainable model parameters. The mapping $f(\cdot)$ therefore learns the
joint relationship between historical trading features, valuation
indicators, external stock relations, and the four-dimensional OHLC
prediction target.
In much of the existing literature, stock price prediction is primarily driven by historical trading data, possibly supplemented by a limited set of fundamental indicators. Building upon these foundations, the proposed \textcolor{blue}{MGAT} further incorporates heterogeneous qualitative signals  to enrich the representation of market dynamics. Moreover, rather than modelling each stock in isolation, \textcolor{blue}{MGAT} explicitly introduces external relational information to capture interdependencies among stocks; such multifaceted relations are re-encoded to better reflect structural and behavioural links among listed companies.
Accordingly, we extend \eqref{eq:multigcn} and describe \textcolor{blue}{MGAT} as:
\begin{equation}
\hat{Y} = f\!\left(
\left[\mathbf{x}_{D-t}, \mathbf{x}_{D-t-1}, \ldots, \mathbf{x}_{D-1},
\mathbf{o}^{*}_{D-t}, \mathbf{o}^{*}_{D-t-1}, \ldots, \mathbf{o}^{*}_{D-1}\right],
\mathcal{G}, \Theta
\right),
\label{eq:mgt}
\end{equation}
where, in addition to the variables defined in Eq.~\eqref{eq:multigcn},
$o^{*}_{D-k}$ denotes the enriched set of valuation and fundamental
indicators, including PE, PB, PS, and PCF, and $\mathcal{G}$ denotes
the set of external relation graphs encoding inter-stock dependencies.
{\color{blue}Equation~\eqref{eq:mgt} states the general form. In the implementation used in this study, MGAT forecasts, at the close of day $t$, the high, low, and closing prices of days $t+1$ and $t+2$ for all stocks from a window of the $T=20$ most recent trading days,
\[
\hat{\mathbf{Y}}_{t+k}=f\!\left(\mathbf{x}_{t-T+1:t},\,\mathbf{o}^{*}_{t-T+1:t},\,\mathcal{G},\,\Theta\right),\qquad k=1,2,
\]
where $\mathbf{x}_{s}$ contains, for each stock, the daily log return $\log(p_{s,i}/p_{s-1,i})$, the logarithms of the high and low prices relative to the close, and the logarithm of the trading value relative to its average over the previous 20 days, and $\mathbf{o}^{*}_{s}$ contains the valuation yields $1/\mathrm{PE}$, $1/\mathrm{PB}$, $1/\mathrm{PS}$, and $1/\mathrm{PCF}$ of day $s$ (yields remain finite when earnings or cash flows change sign). All inputs are observed at the close of day $s\le t$. The targets are the log ratios $\log(p^{c}_{t+k,i}/p_{t,i})$ for $c\in\{\mathrm{high},\mathrm{low},\mathrm{close}\}$, from which price forecasts are recovered as $p_{t,i}\exp(\cdot)$; the forecast high and low are bounded by the forecast close so that each forecast day is a consistent price bar. The opening price is not forecast, because the RL state does not use it and the healthcare data do not contain it.\par}
For presentation convenience, we first introduce the basic notations commonly used in describing GNNs.

A graph can be specified as $\mathcal{G}=(\mathcal{V},\mathcal{E},\mathbf{A})$ \citep{Sun2025GTC}, where $\mathcal{V}$ denotes a set of $\nu$ nodes, $v_i\in\mathcal{V}$, $i=1,2,\ldots,\nu$, and $\mathbf{x}_i$ is the original feature vector associated with node $v_i$.
$\mathcal{E}$ denotes the set of edges indicating connections between nodes, while $A$ is the adjacency matrix encodes the strength of relationships between node pairs, typically constructed according to neighbourhood information.

Let $\mathbf{a}_i^{(k)}$ be the aggregated neighbourhood message for the node $v_i$ at the $k$-th layer, and let $\mathbf{H}_i^{(k)}$ be the representation of the node $v_i$ at the $k$-th layer, with the initial representation $\mathbf{H}_i^{(0)}=\mathbf{x}_i$.
The $k$-th iteration of feature aggregation can be written in an aggregate combined form as
\begin{equation}
\mathbf{a}_i^{(k)}=\mathrm{Aggregate}^{(k)}\!\left(\left\{\mathbf{H}_j^{(k-1)}: j\in \mathcal{N}(v_i)\right\}\right),
\label{eq:gnn_aggregate}
\end{equation}
\begin{equation}
\mathbf{H}_i^{(k)}=\mathrm{Combine}^{(k)}\!\left(\mathbf{a}_i^{(k)}, \mathbf{H}_i^{(k-1)}\right),
\label{eq:gnn_combine}
\end{equation}
where $\mathcal{N}(v_i)$ denotes the neighbourhood set of node $v_i$.
$\mathrm{Aggregate}^{(k)}(\cdot)$ and $\mathrm{Combine}^{(k)}(\cdot)$ denote, respectively, the aggregation operator and the combination operator at layer $k$.
In the following, we instantiate \eqref{eq:gnn_aggregate}--\eqref{eq:gnn_combine} to derive the specific formulations used in our framework.

\subsection{Heterogeneous relation graphs}\label{Heterogeneous relation graphs}

Building on recent graph-learning and graph-transformer advances, including adaptive-filter graph convolution (GCNN) \citep{Andrea2023AdaptiveFiltersGCNN}, structural encoding in Graphormer-style architectures \citep{Ying2021NeurIPSTransformersGraph}, and structure-aware graph Transformers \citep{Chen2022StructureAwareTransformerGraphs}, we design a set of GNN-type relational encoders to extract heterogeneous cross-influence patterns among stocks. In particular, we model the three heterogeneous graphs: (i) mutual shareholding relations, (ii)  board-secretary answering investor inquiry similarity (board-secretary investor relations Q\&A), and (iii) daily \textcolor{blue}{trading-value} co-movement relations. The resulting graphs share the listed companies' set but differ in edge semantics and adjacency construction.

\paragraph{(i) Mutual shareholding relationship graph}
We construct a shareholder relationship network for stocks that share common major shareholders.
Let
\begin{equation}
\mathcal{G}_{\mathrm{SH}}=(\mathcal{V},\mathcal{E}_{\mathrm{SH}},\mathbf{A}_{\mathrm{SH}}),
\end{equation}
where $\mathcal{V}$ is the set of $N$ listed companies and $\mathcal{E}_{\mathrm{SH}}$ denotes the set of shareholding edges.
{\color{blue}
For firm $i$, let $\mathcal{H}_i$ denote the shareholder names in its top-ten shareholder records, and let $h_{ik}$ be the percentage of its shares held by shareholder $k$; a record may list fewer than ten shareholders. The weight of the edge between two distinct firms depends only on the shareholders that appear in both lists:
\begin{equation}
a_{\mathrm{SH},ij}
=
\sum_{k\in\mathcal{H}_i\cap\mathcal{H}_j}
\frac{h_{ik}+h_{jk}}{200},\qquad i\neq j.
\label{eq:ash_def}
\end{equation}
Each common shareholder contributes the mean of its two holding proportions, expressed as a fraction, and firms without a common shareholder receive no edge. The matrix is symmetric, with a zero diagonal except in the case noted below. Shareholder names are matched exactly after removing leading and trailing whitespace, and records are assigned to firms by short name. In the archived banking file, records are available for 13 of the 15 banks (none for PSBC and CZB), and the record of China Merchants Bank (CMB) carries the short name of China Minsheng Bank (CMBC). CMB's list is therefore attributed to CMBC: CMB, PSBC, and CZB are isolated nodes, and the shareholders common to the two lists give CMBC a diagonal entry of 0.504. All experiments use the matrix as archived, so that the reported results can be reproduced exactly (Section~\ref{sec:limitations}). The graph is static and is used without row normalisation.
\par}
Accordingly,
\begin{equation}
\mathbf{A}_{\mathrm{SH}}=
\begin{bmatrix}
a_{\mathrm{SH},11} & a_{\mathrm{SH},12} & \cdots & a_{\mathrm{SH},1N}\\
a_{\mathrm{SH},21} & a_{\mathrm{SH},22} & \cdots & a_{\mathrm{SH},2N}\\
\vdots & \vdots & \ddots & \vdots\\
a_{\mathrm{SH},N1} & a_{\mathrm{SH},N2} & \cdots & a_{\mathrm{SH},NN}
\end{bmatrix}.
\label{eq:ash_mat}
\end{equation}

\paragraph{(ii) Board-secretary investor-relations Q\&A co-mention graph}
{\color{blue}To capture which firms investors discuss together in interactive disclosures, we define a graph from co-mentions in investor questions answered by board secretaries.
Let
\begin{equation}
\mathcal{G}_{\mathrm{QA}}=(\mathcal{V},\mathcal{E}_{\mathrm{QA}},\mathbf{A}_{\mathrm{QA}}).
\end{equation}
Let $\mathcal{Q}$ be the collected corpus of Q\&A records. For each record $q\in\mathcal{Q}$, let $\mathcal{M}(q)\subseteq\mathcal{V}$ be the set of sector firms whose official short name appears verbatim in the question text (exact substring matching; answers are not used). The edge weight is the number of questions that mention both firms:
\begin{equation}
a_{\mathrm{QA},ij}=\sum_{q\in\mathcal{Q}}\mathbb{1}\left[\{i,j\}\subseteq\mathcal{M}(q)\right],\qquad i\neq j,
\label{eq:aqa_def}
\end{equation}
and $a_{\mathrm{QA},ii}=0$. The matrix is symmetric and is computed once over the whole corpus, so the graph is static. Firms that are never co-mentioned with another sector firm are isolated nodes; no imputation is applied to sparse records. This count measures co-attention, i.e.\ how often investors refer to two firms in the same question; it does not measure semantic similarity, sentiment, or the content of the answers.\par}

\paragraph{(iii) Daily trading-value Kendall-correlation graph}
{\color{blue}We further construct a trading-activity co-movement graph from the Kendall rank correlation of daily trading value, i.e.\ the amount traded in CNY:
\begin{equation}
\mathcal{G}_{\mathrm{KD}}=(\mathcal{V},\mathcal{E}_{\mathrm{KD}},\mathbf{A}_{\mathrm{KD}}).
\end{equation}
For firms $i$ and $j$ with daily trading-value series $\{x_\ell\}$ and $\{y_\ell\}$ over a window of $L$ trading days\footnote{We use $L$ rather than $N$ to avoid confusion with the number of listed companies.}, let $n_c$ and $n_d$ be the numbers of concordant and discordant pairs of days, $n_0=L(L-1)/2$, and let $n_1$ and $n_2$ be the numbers of pairs tied in $x$ and in $y$, respectively. We use the tie-adjusted Kendall $\tau_b$,
\begin{equation}
a_{\mathrm{KD},ij}=\frac{n_c-n_d}{\sqrt{(n_0-n_1)(n_0-n_2)}},
\label{eq:kendall_def}
\end{equation}
as implemented in the \texttt{pandas} method \texttt{DataFrame.corr} with \texttt{method="kendall"}. The healthcare series have no missing values. In the banking sector, the trading values of PSBC and CZB are missing before their A-share listings (471 and 461 trading days); these missing values are replaced by the mean of the stock's series over the window before estimation, and the resulting ties are handled by $\tau_b$. The relation matrices used in all experiments are obtained exactly by this procedure. Negative correlations are retained as signed edge weights, and the diagonal equals one. The graph is computed once from a single static window of daily data from January 2, 2018 to December 30, 2022. This window ends before the 2024 selection and 2025 evaluation periods of the RL experiments but overlaps the RL training period and the test years of the first two forecasting folds (Section~\ref{subsec:splits}). The input files label this series ``volume''; it is the daily trading value, not the number of shares traded.\par}

\subsection{Multi-relational graph attention (MGAT)}\label{subsec3}

{\color{blue}MGAT first encodes cross-stock information with graph attention \citep{Velickovic2018GAT} over the relation graphs of Section~\ref{Heterogeneous relation graphs}. Because the three graphs differ in edge semantics and numerical scale, they are combined into one edge set rather than into one weighted operator. Let $\mathbf{A}_r$, $r\in\{\mathrm{SH},\mathrm{QA},\mathrm{KD}\}$, denote the static relation matrices and define
\begin{equation}
\bar{\mathbf{A}}=\sum_{r\in\{\mathrm{SH},\mathrm{QA},\mathrm{KD}\}}\omega_r\,\frac{\max(\mathbf{A}_r,0)}{\max_{i\neq j}\,a_{r,ij}},\qquad \bar a_{ii}=0,
\label{eq:Afused_re}
\end{equation}
where the maximum is taken elementwise, so negative Kendall correlations are excluded, and $\omega_r\in\{0,1\}$ activates relation $r$ ($\omega_r=1$ for all relations in MGAT). The neighbourhood of stock $i$ is
\begin{equation}
\mathcal{N}(i)=\{\,j:\bar a_{ij}>0\,\}\cup\{i\}.
\label{eq:neigh}
\end{equation}
The union of the three graphs covers 92\% of the ordered stock pairs in banking and 82\% in healthcare (minimum degree 6 and 4, respectively), so the relation graphs mainly determine which pairs are excluded, and attention determines the weight of each admissible neighbour. In both sectors, every shareholder edge is also a positive Kendall edge, and the Q\&A layer adds four ordered pairs in banking and none in healthcare; the edge set is therefore determined almost entirely by the trading-value graph, and the qualitative relations can influence the forecasts only to the extent that attention learns to favour the pairs they link (Section~\ref{subsec:posthoc}).

Let $\mathbf{h}^{(k)}_{s,i}$ denote the representation of stock $i$ on day $s$ of the input window after $k$ layers, with $\mathbf{h}^{(0)}_{s,i}=\mathbf{x}_{s,i}\in\mathbb{R}^{F}$. Each graph-attention layer computes
\begin{align}
e^{(k)}_{s,ij}&=\mathrm{LeakyReLU}_{0.2}\!\left(\mathbf{a}_{k,1}^{\top}\mathbf{W}_k\mathbf{h}^{(k)}_{s,i}+\mathbf{a}_{k,2}^{\top}\mathbf{W}_k\mathbf{h}^{(k)}_{s,j}\right),\label{eq:gat_e}\\
\alpha^{(k)}_{s,ij}&=\frac{\exp\big(e^{(k)}_{s,ij}\big)}{\sum_{l\in\mathcal{N}(i)}\exp\big(e^{(k)}_{s,il}\big)},\qquad j\in\mathcal{N}(i),\label{eq:gat_alpha}\\
\mathbf{h}^{(k+1)}_{s,i}&=\mathrm{ELU}\!\left(\sum_{j\in\mathcal{N}(i)}\alpha^{(k)}_{s,ij}\,\mathbf{W}_k\mathbf{h}^{(k)}_{s,j}+\mathbf{b}_k\right),\label{eq:gat_h}
\end{align}
where $\mathbf{W}_k\in\mathbb{R}^{F_{k+1}\times F_k}$, $\mathbf{a}_{k,1},\mathbf{a}_{k,2}\in\mathbb{R}^{F_{k+1}}$, and $\mathbf{b}_k\in\mathbb{R}^{F_{k+1}}$ are trainable and shared across days and stocks. Equations~\eqref{eq:gat_e}--\eqref{eq:gat_h} instantiate the aggregate and combine operators of Eqs.~\eqref{eq:gnn_aggregate}--\eqref{eq:gnn_combine}: the attention coefficients $\alpha^{(k)}_{s,ij}$ depend on the features of both stocks on day $s$, so the weight of a related firm varies with market conditions, whereas a graph convolution applies the same weight on every day. MGAT uses $K=3$ single-head layers of widths 16, 32, and 64 and applies them to each day of the window, giving
\[
\mathbf{H}^{\mathrm{GAT}}_{s}=\big[\mathbf{h}^{(3)}_{s,1},\ldots,\mathbf{h}^{(3)}_{s,N}\big]^{\top}\in\mathbb{R}^{N\times C},\qquad C=64,\qquad s=t-T+1,\ldots,t .
\]

Three variants isolate the role of the relation graphs while keeping the remainder of the forecaster and its training unchanged: a graph-convolution encoder with the fixed symmetric-normalised operator $\mathbf{D}^{-1/2}(\bar{\mathbf{A}}+\mathbf{I})\mathbf{D}^{-1/2}$ and the same layer widths (GCN); an encoder with the same layers applied to each stock separately, which preserves the capacity of the graph branch without cross-stock propagation (no-cross-stock, NCS); and a Transformer without a graph branch (graph-free). The multi-relation encoder of the original submission, which propagated features with a learnable dense matrix $\mathbf{W}_M$ plus learnably weighted relation matrices, is retained as a fourth comparator (MGNN); MGAT contains no free cross-stock matrix, so cross-stock information can flow only along the relation edges.\par}

\subsection{Temporal Transformer for forecasting}\label{subsec4}

The \textcolor{blue}{Transformer encoder \citep{Vaswani2017Attention}} block in MGAT is designed to model temporal dependencies in stock prices while being conditioned on the relational signals extracted by \textcolor{blue}{the graph-attention layers}. 
Different from recurrent encoders that aggregate information sequentially, the Transformer treats historical observations within the input window as tokens and learns content-adaptive interactions across time via self-attention. 
\textcolor{blue}{This general formulation is given below; the implemented configuration, with a window of 20 trading days, is described at the end of this section.}

\paragraph{Transformer-based feature encoding}
Let the raw feature tensor be $\mathbf{X}\in\mathbb{R}^{T\times N\times F}$ and the \textcolor{blue}{graph-attention feature tensor} be $\mathbf{X}_{\mathrm{GNN}}=\mathbf{H}^{\mathrm{GAT}}\in\mathbb{R}^{T\times N\times C}$. 
We first concatenate them along the feature dimension to form the augmented input:
\begin{equation}
\tilde{\mathbf{X}}_{s}
=
\mathrm{Concat}\!\left(\mathbf{X}_{s},\mathbf{X}_{\mathrm{GNN},s}\right)
\in\mathbb{R}^{N\times(F+C)},
\mathbf s
=
\mathbf 1,\ldots,\mathbf T\
\label{eq:concat_re}
\end{equation}

In the subsequent transformer encoder, the $T$ time steps are treated as tokens. Concretely, for each stock $i$, the token sequence is $\tilde{\mathbf{X}}_{i,:}\in\mathbb{R}^{T\times(F+C)}$.

\paragraph{Temporal self-attention computation}
Self-attention computes time-dependent interactions by projecting the token sequence into queries, keys, and values:
\begin{equation}
\mathbf{Q}=\tilde{\mathbf{X}}_{i}\mathbf{W}_{Q},\qquad
\mathbf{K}=\tilde{\mathbf{X}}_{i}\mathbf{W}_{K},\qquad
\mathbf{V}=\tilde{\mathbf{X}}_{i}\mathbf{W}_{V},
\label{eq:qkv}
\end{equation}
where $
\tilde{\mathbf X}_{i}
=
\big[\tilde{\mathbf X}_{t-T+1}(i,:),\ldots,\tilde{\mathbf X}_{t}(i,:)\big]
\in\mathbb R^{T\times(F+C)}
,\mathbf{W}_{Q},\mathbf{W}_{K},\mathbf{W}_{V}\in\mathbb{R}^{(F+C)\times d}$ are trainable projection matrices and $d$ denotes the attention dimension. 
For each stock $i$, the projections are computed $\tilde{\mathbf{X}}_{i}\in\mathbb{R}^{T\times(F+C)}$ so that
$\mathbf{Q}_{i},\mathbf{K}_{i},\mathbf{V}_{i}\in\mathbb{R}^{T\times d}$, and the resulting attention matrix
$\mathrm{softmax}(\mathbf{Q}_{i}\mathbf{K}_{i}^{\top}/\sqrt{d})\in\mathbb{R}^{T\times T}$ explicitly captures the dependency between any pair of time steps. 
The core scaled dot-product attention is:
\begin{equation}
\mathbf{X}_{\mathrm{ATTN}}
=
\mathrm{softmax}\!\left(\frac{\mathbf{Q}\mathbf{K}^{\top}}{\sqrt{d}}\right)\mathbf{V},
\label{eq:attn}
\end{equation}
which is followed by a residual connection and layer normalisation:
\begin{equation}
\mathbf{X}_{L1}
=
\mathrm{LayerNorm}\!\left(\tilde{\mathbf{X}}+\mathbf{X}_{\mathrm{ATTN}}\right)
=
\gamma\frac{\tilde{\mathbf{X}}+\mathbf{X}_{\mathrm{ATTN}}-\mu}{\sqrt{\sigma^{2}+\epsilon}}+\beta,
\label{eq:ln1}
\end{equation}
where $
\tilde{\mathbf X}
=
\big[\tilde{\mathbf X}_{1},\ldots,\tilde{\mathbf X}_{t}\big]
$,$\mu$ and $\sigma^{2}$ are computed along the feature dimension, $\epsilon$ is a small constant and $\gamma,\beta$ are learnable scale and shift parameters. 
In practice, we employ multi-head attention to allow the model to attend to diverse temporal patterns at different subspaces, while preserving the above formulation as the basic attention operator.
\paragraph{Feedforward network and output representation}
The position-wise feedforward network (FFN) is a two-layer MLP applied to each time token:
\begin{equation}
\mathbf{X}_{\mathrm{FFN}}
=
\mathbf{W}_{2}\,\mathrm{ReLU}\!\left(\mathbf{W}_{1}\mathbf{X}_{L1}+\mathbf{b}_{1}\right)
+\mathbf{b}_{2},
\label{eq:ffn}
\end{equation}
where $\mathbf{W}_{1},\mathbf{W}_{2},\mathbf{b}_{1},\mathbf{b}_{2}$ are trainable parameters. 
A second residual connection and layer normalisation yield
\begin{equation}
\mathbf{X}_{L2}=\mathrm{LayerNorm}\!\left(\mathbf{X}_{L1}+\mathbf{X}_{\mathrm{FFN}}\right).
\label{eq:ln2}
\end{equation}
Concretely, we reshape the fused tensor into $N$ independent temporal token sequences and apply self-attention along the temporal axis, yielding a
$T \times T$ attention matrix for each stock, so the transformer models temporal dependency rather than cross-stock attention.
The attention FFN block can be stacked for multiple layers to form the transformer encoder.

\paragraph{Temporal pooling for forecasting and strategy learning}
To obtain a compact representation for each stock at the end of the input window, we apply last-step pooling by selecting the output token at time $T$:
\begin{equation}
\mathbf{X}_{\mathrm{Trf}}
=
\mathbf{X}_{L2}[:,T,:]
\in\mathbb{R}^{N\times d},
\label{eq:last_step_re}
\end{equation}
This vector summarizes the information in the input window and is mapped by an output layer to the \textcolor{blue}{price forecasts} of each stock\textcolor{blue}{; the policy learner receives the forecasts, not this representation}.

\paragraph{Implemented configuration}
{\color{blue}MGAT has $F=8$ input features per stock and day, three graph-attention layers of widths 16, 32, and 64 ($C=64$), and one encoder block (Eqs.~\eqref{eq:qkv}--\eqref{eq:ln2}) applied separately to each stock along the $T=20$ days of the window, with learned positional embeddings, four attention heads, a feedforward width of 128, and dropout 0.1. A linear output layer maps the last-day representation (Eq.~\eqref{eq:last_step_re}) of each stock to six outputs, the high, low, and closing log ratios for $t+1$ and $t+2$; the model has 44{,}846 trainable parameters (GCN and NCS: 44{,}622; MGNN: 44{,}850; graph-free Transformer: 2{,}718). Inputs and targets are standardised with the means and standard deviations of the training rows of each fold, and standardised inputs are clipped to $[-5,5]$. Temporal self-attention operates within each stock, and cross-stock information enters only through the graph-attention representation, which acts as inter-stock attention restricted to the relation edges.\par}

\subsection{Ranking-aware training objective}\label{subsec:objective}
{\color{blue}A forecast made at the close of day $t$ can be acted upon at the close of day $t+1$ at the earliest (Section~\ref{subsec:env}). The first return that such a decision earns is
\begin{equation}
q_{t,i}=\log\frac{C_{t+2,i}}{C_{t+1,i}}=y^{c}_{t,i,2}-y^{c}_{t,i,1},\qquad \hat q_{t,i}=\hat y^{c}_{t,i,2}-\hat y^{c}_{t,i,1},
\label{eq:q}
\end{equation}
where $y^{c}_{t,i,k}=\log(C_{t+k,i}/C_{t,i})$. For allocation across stocks, the cross-sectional ordering of $\hat q_{t,i}$ matters more than its level. The squared-error loss constrains $\hat q$ only indirectly, as the difference of two closing-price outputs whose errors are much larger than the variation of $q$ across stocks. MGAT is therefore trained with
\begin{equation}
\mathcal{L}=\frac{1}{|\mathcal{M}|}\sum_{(t,i,c)\in\mathcal{M}}\big(\hat y_{t,i,c}-y_{t,i,c}\big)^2+\lambda\Big(1-\frac{1}{|\mathcal{T}|}\sum_{t\in\mathcal{T}}\rho_t(\hat{\mathbf{q}}_t,\mathbf{q}_t)\Big),
\label{eq:loss}
\end{equation}
where the first term is the mean squared error over the six standardised targets on the set $\mathcal{M}$ of observed stock-days, $\rho_t$ is the Pearson correlation across the stocks with observed $q_{t,i}$ on day $t$, $\mathcal{T}$ contains the days in the mini-batch with at least five such stocks, and $\lambda=1$. The value of $\lambda$ was fixed before training and was not tuned. All neural forecasters in this study, including the comparators of Section~\ref{subsec3}, are trained with this objective. The linear comparator is a ridge regression of $q$ on the flattened 20-day window of the eight features, pooled over stocks, with the penalty chosen by the validation rank correlation of $q$; a second ridge regression on the same window forecasts the log range $\log(H_{t+1}/L_{t+1})$, with the penalty chosen by the validation mean squared error. Optimisation uses Adam with learning rate $10^{-3}$, batch size 32, at most 100 epochs, and early stopping with patience 10 on the validation loss. Each forecaster is trained with seeds 1--3, and the mean of the three forecasts is used downstream.\par}

\paragraph{Selection of the forecaster}
{\color{blue}The graph encoder was chosen before the evaluation year was examined, on the rank information coefficient of $\hat q$ (IC$_2$, the mean daily Spearman correlation between $\hat q_{t,i}$ and $q_{t,i}$) in the development years 2022 and 2023 (Section~\ref{subsec:splits}). Among the encoders that propagate information across stocks, graph attention had the highest development IC$_2$ in both sectors (banking 0.057, $t=3.5$; healthcare 0.018, $t=1.4$), compared with 0.054 and 0.010 for GCN and $-0.011$ and 0.012 for MGNN. The graph-free Transformer had a higher development IC$_2$ in banking (0.070), and it is therefore retained as a registered comparator, so that any contribution of the relation graphs is tested rather than assumed (Section~\ref{subsec:hypotheses}).\par}

\section{Reinforcement learning for investment strategy}\label{sec4}

{\color{blue}This section describes the decision layer of MGAT-FinRL and its interface with the forecaster. The aim is to measure how much the forecasts change the allocations learned by standard DRL agents when everything else in the environment is held fixed.\par}

\subsection{FinRL-based portfolio optimization}\label{FinRL-based portfolio optimization}

{\color{blue}As illustrated in Fig.~\ref{fig:finrl}, MGAT is trained first and then held fixed; its out-of-sample forecasts are part of the agent's state at every decision in the training, selection, and evaluation periods.\par}
\begin{figure*}[!t]
    \centering
    \includegraphics[width=\textwidth]{fig2_finrl.pdf}
    \caption{{\color{blue}FinRL layer of MGAT-FinRL. The state combines the 252-day covariance matrix and four technical indicators computed from observed prices with two forecast rows produced by MGAT at the close of the decision day. The agent's action is projected onto long-only weights, which are executed at the next close in a share-and-cash ledger with trading costs. The lower bar shows the walk-forward use of the data.}}
    \label{fig:finrl}
\end{figure*}

We evaluate the decision-making layer under the FinRL framework with a compact set of widely used continuous-control algorithms: PPO\citep{schulman2017ppo}, A2C\citep{mnih2016a3c}, SAC\citep{haarnoja2018sac}, TD3\citep{fujimoto2018td3}, and DDPG\citep{lillicrap2016ddpg}. They span the main design axes in practical portfolio control. PPO and A2C are on-policy actor-critic methods with stable updates under non-stationary return distributions; PPO further constrains policy updates by clipping and can use entropy regularisation to encourage exploration. DDPG and TD3 are off-policy deterministic actor-critic methods for continuous control; TD3 mitigates Q-value overestimation through target policy smoothing and delayed actor updates. SAC is an off-policy stochastic method that maximises expected return plus an entropy term, which promotes exploration. All agents use the Stable-Baselines3 implementations with the FinRL portfolio-allocation interface \citep{Liu2021FinRL}; the core algorithms are not modified.

\paragraph{Coupling with MGAT}
{\color{blue}At the close of day $t$, MGAT forecasts $\hat q_{t,i}$ (Eq.~\eqref{eq:q}) and the high and low of day $t+1$ for every stock. Two rows are added to the state: the cross-sectional $z$-score of the predicted tradable-horizon return,
\begin{equation}
z_{t,i}=\frac{\hat q_{t,i}-\frac1N\sum_{j}\hat q_{t,j}}{\mathrm{sd}_j(\hat q_{t,j})},
\label{eq:zq}
\end{equation}
and the predicted log range $\hat r_{t,i}=\hat y^{h}_{t,i,1}-\hat y^{l}_{t,i,1}$, standardised with its mean and standard deviation over the RL training period. Each forecast is out of sample: it is produced by the forecaster of the walk-forward fold whose test year contains day $t+1$ (Section~\ref{subsec:splits}). Because the forecaster is held fixed, portfolio rewards do not update MGAT, and all agents and seeds receive identical forecasts.

The same environment is run with forecasts from four sources, each producing the same two rows: MGAT; the graph-free Transformer trained with the same objective; the ridge regressions of Section~\ref{subsec:objective}; and a naive source with $\hat q_{t,i}=0$ (so $z_{t,i}=0$) and the observed range $\log(H_{t,i}/L_{t,i})$ of day $t$ in place of the predicted range. The naive source fixes the state dimension and removes forecast information, so differences between the MGAT and naive agents measure the contribution of the forecasts to the learned allocations.\par}

\paragraph{State space}
{\color{blue}The state at decision time $t$ is
\begin{equation}
s_t=\left[\boldsymbol{\Sigma}_{t};\ \mathbf{D}_{t};\ \mathbf{z}_{t}^{\top};\ \tilde{\mathbf{r}}_{t}^{\top}\right]\in\mathbb{R}^{(N+6)\times N},
\label{eq:state}
\end{equation}
where $\boldsymbol{\Sigma}_{t}\in\mathbb{R}^{N\times N}$ is the sample covariance matrix of the last 252 daily close-to-close returns, $\mathbf{D}_{t}\in\mathbb{R}^{4\times N}$ stacks MACD, the 30-day RSI, the 30-day CCI, and the 30-day DX of each asset computed with the \texttt{stockstats} library from observed daily high, low, and close prices up to day $t$, $\mathbf{z}_t$ is given by Eq.~\eqref{eq:zq}, and $\tilde{\mathbf{r}}_t$ is the standardised predicted range. For $N=15$, the state has $21\times15$ entries for every forecast source.\par}

\paragraph{Action space and feasibility constraints}
At each decision time $t$, the policy outputs a real-valued vector $\mathbf{c}_{t}\in[-1,1]^{N}$, which is mapped to a long-only portfolio weight vector using a softmax projection:
\begin{equation}
w_{t,i}=\frac{\exp(c_{t,i})}{\sum_{j=1}^{N}\exp(c_{t,j})},\qquad
\mathbf{w}_{t}=(w_{t,1},\ldots,w_{t,N})\in\Delta^{N},
\label{eq:portfolio_softmax}
\end{equation}
where $\Delta^{N}=\{\mathbf{w}\in\mathbb{R}^{N}\mid w_i\ge 0,\ \sum_{i=1}^{N}w_i=1\}$. {\color{blue}With actions bounded in $[-1,1]$, the largest weight ratio between two stocks is $e^{2}\approx7.4$ and a single weight lies between 0.010 and 0.35, so the policies tilt an approximately diversified portfolio rather than concentrate it.} All agents operate on the same continuous action space.

\paragraph{Reward and training objective}
{\color{blue}
At decision time $t$ the agent chooses target weights $\mathbf{w}_t$ using information available at the close of day $t$, and the trade is executed at the close of day $t+1$ (Section~\ref{subsec:env}). Let $V_t$ denote the account value at the close of day $t$, after any trade executed on that day. The immediate reward is the scaled net one-step portfolio return,
\begin{equation}
R_{t+1}=\kappa_R\,\frac{V_{t+1}-V_t}{V_t},\qquad \kappa_R=0.1,
\label{eq:rl_reward}
\end{equation}
where $V_{t+1}$ is computed from observed prices and is net of the trading costs charged on day $t+1$ (Eq.~\eqref{eq:trading_cost}). The agents maximise the expected discounted sum of rewards with discount factor $\gamma=0.99$. Each training episode covers the RL training period once, starting from an all-cash account of $V_0=10^6$.

A common alternative in FinRL-style environments is to return the updated account value, $R_{t+1}=V_{t+1}$, as the immediate reward. Wealth-level rewards are non-stationary in scale and weight late, high-wealth steps more heavily, whereas the one-step return in Eq.~\eqref{eq:rl_reward} has a stable scale across the episode and lets trading costs enter the learning signal directly. Reward variants that penalise volatility or drawdown were not evaluated.
\par}

\section{Experimental Settings and Performance Analysis}\label{sec5}

\subsection{Datasets and backtesting protocol}
To evaluate the proposed MGAT-FinRL framework, we conduct backtesting on two equity sectors that represent markedly different market structures and information environments. The first sector consists of 15 publicly listed Chinese banks, a relatively stable and regulation-sensitive sector. The second sector contains 15 healthcare innovation stocks listed on the Shanghai and Shenzhen stock exchanges, a growth-oriented sector with higher idiosyncratic volatility. The two-sector design allows a comparison of model behaviour across distinct sectoral regimes.

\paragraph{Data sources and features}
{\color{blue}For each stock, we collect daily high, low, and closing prices, the daily trading value (the amount traded in CNY), and the valuation ratios PE (trailing twelve months), PB, PS, and PCF from January 2018 to December 2025. Prices are unadjusted, and the ledger neither credits cash dividends nor adjusts share quantities for corporate actions; the reported performance therefore excludes dividend income (Section~\ref{sec:limitations}).

The 2018--2024 panels are those of the original study. The 2025 observations were added for this revision. Daily prices and trading values were downloaded from Tushare (\texttt{daily}). The 2025 valuation ratios were rebuilt to match the definitions of the existing panels, which were identified on the 2024 overlap: PE and PB from the daily valuation records of the Eastmoney data centre, and PS and PCF as total market value divided by, respectively, total operating income and the net increase in cash and cash equivalents of the previous fiscal year, from the annual statements in Eastmoney F10. Before any model used a 2025 observation, the downloads were checked against the existing panels over all 242 trading days of 2024: high, low, and closing prices and trading values agree exactly on all 3{,}630 stock-days in each sector; the rebuilt PE, PS, and PCF agree with a median absolute relative difference below $10^{-4}$ and PB with a median difference of 0.3\%, with identical signs throughout. The extended panels contain the 2024 rows unchanged and 243 trading days of 2025.\par}

In addition, we assemble heterogeneous relational information to construct the graphs of Section~\ref{Heterogeneous relation graphs}: top-ten shareholder holdings for the 15 firms of each sector and board-secretary investor-relations Q\&A records.
{\color{blue}
The shareholder data are static snapshots exported from Wind; their record coverage is specified in Section~\ref{Heterogeneous relation graphs}. The Q\&A corpus contains 6{,}021 banking records and 7{,}868 healthcare records. Table~\ref{tab:timeline} documents their source periods, availability, and timestamp limitations.
\par}

\paragraph{Stock universes}
{\color{blue}
The fixed stock lists are provided in Appendix~\ref{app:universe}. The banking universe comprises six large state-owned banks and nine listed national joint-stock commercial banks, all listed by the end of 2019. The healthcare universe was compiled from Wind and covers medical services, diagnostics, eye care, and contract research; it is not an index constituent list. All healthcare stocks were listed before 2018, and all 30 firms remained listed through 2025, which introduces survivorship bias. PSBC and CZB listed on December~10 and November~26,~2019, respectively; their pre-listing placeholder observations are excluded as forecasting targets, with input features set to the training mean, and the banking market history used for RL states starts on December~10,~2019. Stock-days without trading (three in healthcare in 2025) are excluded as forecasting targets and carried at the previous close in the backtest.
\par}

\paragraph{Training, selection, and evaluation periods}
\label{subsec:splits}
{\color{blue}
Forecasting uses four expanding-window folds with test years 2022, 2023, 2024, and 2025. Each fold is trained on origins from 2018 onwards whose two targets fall no later than the end of the second year before the test year, early-stopped on the intervening year, and tested on the origins whose next trading day lies in the test year; standardisation uses training rows only. A forecast made at the close of day $t$ is therefore always produced by a model that was fitted and stopped on data ending before the year of day $t+1$.

The RL periods are aligned with these folds. Agents are trained on January~4, 2022--December~29, 2023 (484 trading days), using the forecasts of folds 2022 and 2023 in the state. For each agent and seed, checkpoints saved at 20\%, 40\%, 60\%, 80\%, and 100\% of the training budget are compared on 2024 (242 days, forecasts of fold 2024), and the checkpoint with the highest net Sharpe ratio under cost setting F1 is retained. The selected policies are frozen and evaluated once on January~2--December~31, 2025 (243 days), with forecasts of fold 2025 and under all cost settings. Each agent is trained with seeds 1--5, and all agents, seeds, and cost settings are reported.

The years play distinct roles. 2022 and 2023 are development years: the forecaster was chosen on them (Section~\ref{subsec:objective}). 2024 was examined in earlier versions of this study and serves here only for checkpoint selection. 2025 was not used in any earlier analysis; its data were downloaded after the protocol of Section~\ref{subsec:hypotheses} had been fixed.
\par}

\paragraph{Pre-registered hypotheses}\label{subsec:hypotheses}
{\color{blue}
The protocol, forecaster, forecast sources, agents, budgets, seeds, selection rule, cost settings, and tests were committed to the public code repository before any 2025 observation was downloaded, together with two dated amendments that concern only the source of the 2025 valuation ratios. The outcome is the annualised net Sharpe ratio in 2025 under F1, compared in pairs with identical agent and seed over the 25 policies of each source (five agents $\times$ five seeds):
\begin{itemize}
\item H1 (primary): the MGAT agents attain a higher Sharpe ratio than the naive agents, in each sector (Holm correction across the two sectors);
\item H2 (secondary): the MGAT agents attain a higher Sharpe ratio than the graph-free Transformer agents (H2a) and than the ridge agents (H2b) (Holm correction within each sector over H1, H2a, and H2b);
\item H3 (forecasting): the 2025 IC$_2$ of MGAT is positive in each sector (one-sided $t$-test on daily coefficients, Holm correction across sectors).
\end{itemize}
Differences in Sharpe ratio are tested with a paired circular block bootstrap of daily returns (block length 10 days, with 5 and 20 as checks; 2{,}000 resamples; the same resampled dates for all policies), using one-sided centred $p$-values at the 5\% level. The GCN, no-cross-stock, and MGNN forecasters enter PPO agents with seeds 1--5 as a descriptive ablation.
\par}

\paragraph{Information timeline}
{\color{blue}
Table~\ref{tab:timeline} summarises data availability and execution timing. All three relation graphs use information that ends by June 2023: the Kendall graph uses trading values from January 2, 2018 to December 30, 2022, the shareholder snapshots date from 2022--2023, and the Q\&A records from 2018--2022. They therefore contain no information from the selection year 2024 or the evaluation year 2025. They do overlap the test years of folds 2022 and 2023, so the forecasts that enter the RL training states are retrospective in this respect (Section~\ref{sec:limitations}).
\par}

\begin{table*}[!t]
\centering
\small
\caption{{\color{blue}Data availability and forecasting--trading timeline.}}
\label{tab:timeline}
\setlength{\tabcolsep}{3pt}
\renewcommand{\arraystretch}{1.12}
\color{blue}
\begin{tabular}{@{}p{2.3cm}p{3.3cm}p{3.3cm}p{3.5cm}@{}}
\toprule
Input or stage & Source period & Availability or protocol cutoff & Use \\
\midrule
Daily market data & Trading day $t$ (2018--2025) & After the close of $t$ & Forecasting features; RL states; accounting \\
Valuation ratios & Trading day $t$ (2018--2025) & After the close of $t$ & Forecasting features only \\
Top-ten shareholders & Healthcare: year-end 2022; banking: lists exported in June 2023 & Healthcare: reports released in 2023; banking: release dates not verified & Static $\mathbf{A}_{\mathrm{SH}}$ \\
Investor-relations Q\&A & Records dated 2018--2022 & Archived record timestamps & Static $\mathbf{A}_{\mathrm{QA}}$ \\
Trading-value Kendall graph & January 2, 2018--December 30, 2022 & After the source window ends & Static $\mathbf{A}_{\mathrm{KD}}$ \\
Forecasting folds & Test years 2022, 2023, 2024, 2025 & Fitted and early-stopped before the test year & 2022--2023: RL training states; 2024: selection states; 2025: evaluation states \\
RL training & January 4, 2022--December 29, 2023 & -- & Five agents, seeds 1--5, cost setting F1 \\
Checkpoint selection & 2024 & After the 2024 close & One checkpoint per agent and seed (F1) \\
Portfolio evaluation & 2025 & State after the close of $t$; execution at the close of $t+1$ & Frozen policies, evaluated once; F0--F6 \\
\bottomrule
\end{tabular}
\end{table*}

{\color{blue}
For the March~3,~2025 forecast origin (a Monday), MGAT forecasts the prices of March~4 and~5 from the 20 trading days ending on March~3. The RL state uses these forecasts and observed data through the March~3 close; the weights chosen after that close are executed at the March~4 close, and the reward of Eq.~\eqref{eq:rl_reward} for March~4 reflects the return on the existing holdings and the execution costs, while the new positions earn the return from the March~4 close onwards, which is the return $q$ ranked by the forecaster.
\par}

\subsection{Relationship graphs}\label{subsec:graphs}

Following the construction principles in Section~\ref{Heterogeneous relation graphs}, we build three relation graphs for each sector:
(i) the mutual shareholding graph $\mathcal{G}_{\mathrm{SH}}$,
(ii) the board-secretary investor-relations Q\&A \textcolor{blue}{co-mention} graph $\mathcal{G}_{\mathrm{QA}}$, and
(iii) the daily \textcolor{blue}{trading-value} Kendall-correlation graph $\mathcal{G}_{\mathrm{KD}}$.
Figures~\ref{fig:fig3} and~\ref{fig:fig4} \textcolor{blue}{visualise} the resulting graphs for the banking sector (left column) and the healthcare innovation sector (right column).

\paragraph{(i) Mutual shareholding ($\mathcal{G}_{\mathrm{SH}}$)}
The banking sector exhibits a denser connectivity pattern, suggesting stronger cross-firm linkages induced by overlapping major shareholders.
By contrast, the healthcare innovation sector contains more disconnected components, indicating that ownership linkages are concentrated on a smaller subset of firms. \textcolor{blue}{In the healthcare graph, six firms share no top-ten shareholder with another sector firm. In the banking graph, CMB, PSBC, and CZB are isolated because of the record attribution described in Section~\ref{Heterogeneous relation graphs}, so the graph understates the ownership links among the 15 banks.}

\paragraph{(ii) Board-secretary investor-relations Q\&A co-mentions ($\mathcal{G}_{\mathrm{QA}}$)}
\textcolor{blue}{In the banking sector, 63 of the 105 firm pairs are co-mentioned at least once, and one bank is never co-mentioned with another sector firm.}
\textcolor{blue}{In the healthcare innovation sector, only seven pairs are co-mentioned and six firms are isolated, so this layer is very sparse; the co-mentions form small local clusters.}

\paragraph{(iii) \textcolor{blue}{Trading-value} Kendall correlation ($\mathcal{G}_{\mathrm{KD}}$)}
The banking sector presents stronger co-movement in \textcolor{blue}{trading-value} dynamics, which is consistent with higher sector-level synchronization in trading activity. The healthcare innovation sector exhibits weaker and more dispersed Kendall correlations, suggesting more stock-specific trading-activity fluctuations and less sector-wide synchronization.
\begin{figure*}[t]
    \centering

    % Row 1: (a) and (d)
    \begin{subfigure}[t]{0.48\textwidth}
        \centering
        \includesvg[width=\linewidth]{fig4(a)}
        \caption{Banking sector: mutual shareholding graph.}
    \end{subfigure}\hfill
    \begin{subfigure}[t]{0.48\textwidth}
        \centering
        \includesvg[width=\linewidth]{fig4(d)}
        \caption{Healthcare innovation sector: mutual shareholding graph.}
    \end{subfigure}

    \vspace{0mm}

    % Row 2: (b) and (e)
    \begin{subfigure}[t]{0.48\textwidth}
        \centering
        \includesvg[width=\linewidth]{fig4(b)}
        \caption{Banking sector: board-secretary investor-relations Q\&A co-mention graph.}
    \end{subfigure}\hfill
    \begin{subfigure}[t]{0.48\textwidth}
        \centering
        \includesvg[width=0.9\linewidth]{fig4(e)}
        \caption{Healthcare innovation sector: board-secretary investor-relations Q\&A co-mention graph.}
    \end{subfigure}
    \caption{Mutual-shareholding and investor-relations Q\&A co-mention
graphs for the banking and healthcare innovation sectors.}
    \label{fig:fig3}
\end{figure*}

\begin{figure*}[t]
    \centering


    % Row 3: (c) and (f)
    \begin{subfigure}[t]{0.48\textwidth}
        \centering
        \includesvg[width=\linewidth]{fig4(c)}
        \caption{Banking sector: \textcolor{blue}{trading-value} Kendall correlation matrix.}
    \end{subfigure}\hfill
    \begin{subfigure}[t]{0.48\textwidth}
        \centering
        \includesvg[width=\linewidth]{fig4(f).svg}
        \caption{Healthcare innovation sector: \textcolor{blue}{trading-value} Kendall correlation matrix.}
    \end{subfigure}

    \caption{\textcolor{blue}{Trading-value} Kendall-correlation matrices for the banking and healthcare innovation sectors.}
    \label{fig:fig4}
\end{figure*}

Overall, different relation types reveal different structural patterns, and their density varies substantially across the two sectors. \textcolor{blue}{Whether these patterns improve downstream allocation is a separate empirical question, examined in Section~\ref{subsec:ablation}.}

\subsection{Forecasting results}\label{sec:price_pred}

{\color{blue}Table~\ref{tab:forecast_eval} reports forecast accuracy in the evaluation year 2025, and Table~\ref{tab:forecast_years} in Appendix~\ref{app:model_complexity} reports all four test years. Each forecast is produced by the fold whose model was fitted and stopped before the test year, and the reported values are means of three seeds. For each price channel $c\in\{\mathrm{high},\mathrm{low},\mathrm{close}\}$ of day $t+1$, accuracy is measured on the log ratios of Section~\ref{subsec:objective}, whose errors equal the errors of the log prices,
\begin{equation}
\mathrm{RMSE}_c=\Big(\frac{1}{|\mathcal{S}_c|}\sum_{(t,i)\in\mathcal{S}_c}\big(\hat y^{c}_{t,i,1}-y^{c}_{t,i,1}\big)^2\Big)^{1/2},
\label{eq:rmse}
\end{equation}
where $\mathcal{S}_c$ contains all forecast origins $t$ of the test year and all stocks $i$ with an observed target, so the average runs over both dates and stocks and each channel is evaluated separately. Tables report $\mathrm{RMSE}_c/\mathrm{RMSE}^{\mathrm{prev}}_c-1$, where the benchmark uses the previous day's values ($\log(H_{t,i}/C_{t,i})$, $\log(L_{t,i}/C_{t,i})$, and 0 for the close). The ranking measure at the tradable horizon is
\begin{equation}
\mathrm{IC}_2=\frac{1}{|\mathcal{T}_{\mathrm{test}}|}\sum_{t\in\mathcal{T}_{\mathrm{test}}}\rho^{S}_t\big(\hat{\mathbf{q}}_t,\mathbf{q}_t\big),
\label{eq:ic2}
\end{equation}
where $\rho^{S}_t$ is the Spearman correlation across the stocks with observed $q_{t,i}$ (Eq.~\eqref{eq:q}) on day $t$, with a $t$-statistic computed from the daily coefficients.

\paragraph{Price levels}
All neural forecasters predict the next day's high and low more accurately than the previous day's values in every test year. In 2025, MGAT reduces the root mean squared error of the log high by 18.9\% in banking and 15.5\% in healthcare, and that of the log low by 18.9\% and 28.5\%; across 2022--2025 the reductions range from 7\% to 29\%. The intraday range is persistent and partly predictable from the recent window, whereas the closing return is not: no forecaster improves on the previous close (changes of $+0.1\%$ to $+6.1\%$ in 2025), consistent with weak-form efficiency at the daily horizon.

\paragraph{Ranking at the tradable horizon}
The registered forecasting hypothesis H3 concerns the rank correlation between $\hat q_t$ and the realised return from the close of $t+1$ to the close of $t+2$. In 2025, the IC$_2$ of MGAT is 0.052 in healthcare ($t=2.28$, one-sided $p=0.012$, Holm-adjusted $p=0.024$) and 0.014 in banking ($t=0.70$, $p=0.24$). H3 is therefore supported in healthcare and not in banking. The healthcare signal is shared by the neural forecasters trained with the ranking-aware objective (IC$_2$ between 0.038 and 0.053, $t$ between 1.8 and 2.4), and the paired daily differences between MGAT and the graph-free, GCN, and no-cross-stock forecasters are small, with 95\% block-bootstrap intervals that include zero (MGAT minus graph-free: 0.005, $[-0.018, 0.027]$). In banking, MGAT ranks the tradable return better than the ridge regression (0.044, $[0.015, 0.074]$), whereas its differences from the other neural forecasters include zero. Across the four test years, the IC$_2$ of MGAT varies between $-0.007$ (banking, 2024) and 0.081 (banking, 2023); its sign is positive in seven of the eight sector-years.

These results indicate that the ranking information comes mainly from the inputs and the training objective, with a modest and statistically unresolved contribution of the relation graphs, and that its strength varies across years. The cross-sectional rank correlations of 0.01--0.08 are of the size commonly reported for daily equity signals; they are informative for ranking but explain a small fraction of the cross-sectional variation in returns.\par}

\begin{table*}[!t]\centering\small\color{blue}
\caption{Forecast accuracy in the evaluation year 2025 (fold 2025: trained on 2018--2023 targets, early-stopped on 2024). High, Low, Close: change in the root mean squared error of the next-day log high, low, and close relative to the previous-day values (\%; negative values are improvements). IC$_2$: mean daily Spearman correlation between the predicted and realised return from the close of $t+1$ to the close of $t+2$, with its $t$-statistic (242 days). Ridge forecasts only this return. Means of three seeds.}
\label{tab:forecast_eval}
\setlength{\tabcolsep}{4pt}
\resizebox{\ifdim\width>\textwidth\textwidth\else\width\fi}{!}{%
\begin{tabular}{lrrrrrrrr}\toprule
 & \multicolumn{4}{c}{Banking} & \multicolumn{4}{c}{Healthcare} \\ \cmidrule(lr){2-5}\cmidrule(lr){6-9}
Forecaster & High & Low & Close & IC$_2$ ($t$) & High & Low & Close & IC$_2$ ($t$) \\ \midrule
MGAT & $-$18.9 & $-$18.9 & +6.1 & 0.014 (0.7) & $-$15.5 & $-$28.5 & +0.5 & 0.052 (2.3) \\
GCN & $-$19.7 & $-$18.1 & +5.4 & 0.025 (1.3) & $-$15.8 & $-$29.9 & +0.4 & 0.053 (2.4) \\
No cross-stock & $-$19.8 & $-$19.8 & +2.5 & 0.032 (1.7) & $-$15.7 & $-$28.2 & +1.2 & 0.044 (2.1) \\
MGNN & $-$14.9 & $-$18.9 & +4.1 & 0.000 (0.0) & $-$14.5 & $-$27.4 & +1.9 & 0.038 (1.8) \\
Graph-free & $-$20.2 & $-$20.4 & +0.2 & 0.027 (1.3) & $-$14.8 & $-$27.4 & +0.1 & 0.047 (2.1) \\
Ridge & -- & -- & -- & $-$0.030 ($-$1.6) & -- & -- & -- & 0.026 (1.2) \\
\bottomrule\end{tabular}}
\end{table*}

\subsection{Trading environment and implementation details}
\label{subsec:env}
{\color{blue}
We implement a portfolio-allocation environment with an explicit share-and-cash ledger and initial capital $V_0=10^6$. State and action definitions follow Section~\ref{FinRL-based portfolio optimization}. Decisions made after the close of day $t$ are executed at the observed close of day $t+1$. Portfolios are long-only and fully invested, with no leverage, fractional shares, and at most one rebalancing per trading day.

Let $B_{t+1}$ and $S_{t+1}$ denote the buy and sell notionals, computed from price-drifted holdings and fee-aware target holdings. Trading costs are
\begin{equation}
C_{t+1}=(c+s)(B_{t+1}+S_{t+1})+\tau_{t+1}S_{t+1},
\label{eq:trading_cost}
\end{equation}
where $c$ is the commission rate and $s$ is a slippage allowance charged as a cash deduction without adjusting execution prices. The sell-side stamp duty $\tau_{t+1}$ is 10 basis points (bp) before August~28,~2023 and 5~bp thereafter. Target holdings are scaled to maintain self-financing:
\begin{equation}
V_{t+1}=V^{-}_{t+1}-C_{t+1},
\qquad
V^{-}_{t+1}=c^{\mathrm{cash}}_t+\sum_{i=1}^{N}q_{t,i}p_{t+1,i},
\label{eq:portfolio_value_update}
\end{equation}
where $q_{t,i}$ and $c^{\mathrm{cash}}_t$ are the holdings and cash after the day-$t$ trade, and $p_{t+1,i}$ is the observed closing price. Initial purchases incur costs; terminal holdings are marked to market without liquidation charges. Gross turnover is the daily average of $(B_{t+1}+S_{t+1})/V^{-}_{t+1}$.

We evaluate seven commission/slippage settings in bp: F0 (0/0, without stamp duty), F1 (3/2), F2 (5/2), F3 (10/2), F4 (20/2), F5 (3/0), and F6 (3/5). F1--F6 include stamp duty. Training and checkpoint selection use F1, and each selected policy is evaluated under all seven settings without retraining. Newly purchased shares cannot be sold before the following trading day; minimum commissions, board lots, price limits, market impact, and liquidity constraints are not modelled.

All agents use Stable-Baselines3 2.4.0 with multilayer-perceptron policies and seeds 1--5. Hyperparameters and metric definitions are provided in Appendices~\ref{app:hyper} and~\ref{app:portfolio_metrics}. Sharpe ratios use a zero risk-free rate. A constant annual rate $r_f$ would lower each Sharpe ratio by approximately $r_f/\sigma_{\mathrm{ann}}$; because the annualised volatilities of the policies lie within 14.9--16.2\% in banking and 27.8--34.5\% in healthcare, a rate of 1.5\% would change the differences between policies of the same sector by at most 0.01. As reference strategies, we report daily equal weighting, equal-weight buy-and-hold, and forecast-ranking rules (Section~\ref{subsec:rules}), all with the same ledger and costs.
\par}

\subsection{Portfolio results in the evaluation year}\label{subsec:comparison}

{\color{blue}Table~\ref{tab:main_f1} summarises the 2025 performance of the 25 policies of each forecast source under F1, and Table~\ref{tab:paired_boot} reports the registered tests.

\paragraph{Pre-registered hypotheses}
The primary hypothesis H1 is not supported in either sector. In healthcare, the MGAT agents attain a mean Sharpe ratio of 1.20, compared with 1.15 for the naive agents (paired difference 0.043, 90\% interval $[-0.066, 0.164]$, $p=0.27$; Holm-adjusted $p=0.55$); the MGAT agents also have the highest mean annualised return (36.8\%) and Calmar ratio (1.96) and the smallest mean drawdown (18.9\%) of the four sources. In banking, the MGAT agents attain 0.54, compared with 0.58 for the naive agents (difference $-0.040$, $[-0.108, 0.023]$, $p=0.83$). The secondary hypotheses H2a and H2b are not supported either: in healthcare, the differences from the graph-free and ridge agents are 0.048 ($p=0.093$) and 0.042 ($p=0.124$), with a Holm-adjusted $p$ of 0.28 for both; in banking, both differences are slightly negative. The conclusions are the same for block lengths of 5 and 20 days.

\paragraph{Size of the effects}
The differences between forecast sources are small relative to the dispersion across agents and seeds. The standard deviation of the Sharpe ratio across the 25 policies of one source is 0.08--0.19, whereas the mean differences between sources are at most 0.05. In healthcare, the MGAT policy has a higher Sharpe ratio than the naive policy with the same agent and seed in 13 of 25 pairs, and in banking in 11 of 25. All RL policies remain close to an equally weighted portfolio, and daily equal weighting attains a Sharpe ratio of 0.61 in banking and 1.23 in healthcare, which 24 of the 100 banking policies and 25 of the 100 healthcare policies exceed. The two sectors also differ in their 2025 market: daily equal weighting returned 8.4\% in banking and 37.6\% in healthcare, with maximum drawdowns of 15.5\% and 19.1\%.\par}

\begin{table*}[!t]\centering\small\color{blue}
\caption{Portfolio performance in the evaluation year 2025 under cost setting F1 (3~bp commission, 2~bp slippage, stamp duty). RL rows: mean\,$\pm$\,standard deviation over the 25 policies of each forecast source (five agents $\times$ five seeds), each trained on 2022--2023 and selected on 2024. AR: annualised return; SR: annualised Sharpe ratio; MDD: maximum drawdown; turnover: mean daily gross turnover. Reference strategies are deterministic.}
\label{tab:main_f1}
\setlength{\tabcolsep}{4pt}
\resizebox{\ifdim\width>\textwidth\textwidth\else\width\fi}{!}{%
\begin{tabular}{llrrrrr}\toprule
Sector & Strategy & AR (\%) & SR & MDD (\%) & Calmar & Turnover (\%/day) \\ \midrule
Banking & Naive agents & 8.1\,$\pm$\,1.8 & 0.58\,$\pm$\,0.11 & 15.9\,$\pm$\,0.8 & 0.51\,$\pm$\,0.12 & 5.8 \\
 & Ridge agents & 7.6\,$\pm$\,2.3 & 0.55\,$\pm$\,0.14 & 15.7\,$\pm$\,1.1 & 0.49\,$\pm$\,0.16 & 7.1 \\
 & Graph-free agents & 7.7\,$\pm$\,1.9 & 0.56\,$\pm$\,0.12 & 15.5\,$\pm$\,0.9 & 0.50\,$\pm$\,0.14 & 4.6 \\
 & MGAT agents & 7.4\,$\pm$\,1.4 & 0.54\,$\pm$\,0.08 & 15.6\,$\pm$\,1.0 & 0.48\,$\pm$\,0.12 & 5.6 \\
 & Equal weight, daily & 8.4 & 0.61 & 15.5 & 0.54 & 1.0 \\
 & Equal weight, buy-and-hold & 8.6 & 0.61 & 15.5 & 0.56 & 0.4 \\
 & Top-3 MGAT rule & $-$10.1 & $-$0.54 & 24.7 & $-$0.41 & 113.5 \\
\midrule
Healthcare & Naive agents & 34.7\,$\pm$\,4.9 & 1.15\,$\pm$\,0.12 & 19.6\,$\pm$\,1.0 & 1.78\,$\pm$\,0.30 & 9.7 \\
 & Ridge agents & 34.8\,$\pm$\,6.4 & 1.15\,$\pm$\,0.17 & 19.1\,$\pm$\,1.1 & 1.83\,$\pm$\,0.40 & 9.4 \\
 & Graph-free agents & 34.8\,$\pm$\,7.5 & 1.15\,$\pm$\,0.19 & 19.4\,$\pm$\,0.8 & 1.80\,$\pm$\,0.41 & 9.7 \\
 & MGAT agents & 36.8\,$\pm$\,7.4 & 1.20\,$\pm$\,0.16 & 18.9\,$\pm$\,1.1 & 1.96\,$\pm$\,0.42 & 8.4 \\
 & Equal weight, daily & 37.6 & 1.23 & 19.1 & 1.97 & 1.8 \\
 & Equal weight, buy-and-hold & 38.0 & 1.19 & 20.4 & 1.86 & 0.4 \\
 & Top-3 MGAT rule & 28.7 & 1.04 & 15.0 & 1.91 & 66.3 \\
\bottomrule\end{tabular}}
\end{table*}
\begin{table*}[!t]\centering\small\color{blue}
\caption{Pre-registered tests in 2025 (F1). $\Delta$SR: mean paired difference in annualised Sharpe ratio over the 25 agent--seed pairs; 90\% interval and one-sided $p$-value from a paired circular block bootstrap of daily returns (block length 10; 2{,}000 resamples); $p_{5}$ and $p_{20}$: block lengths 5 and 20. Holm: H1 across sectors; H1--H2b within each sector.}
\label{tab:paired_boot}
\setlength{\tabcolsep}{3pt}
\resizebox{\ifdim\width>\textwidth\textwidth\else\width\fi}{!}{%
\begin{tabular}{llrrrrrr}\toprule
Sector & Hypothesis & $\Delta$SR & 90\% interval & $p$ & $p_{5}$ & $p_{20}$ & Holm $p$ \\ \midrule
Banking & H1: MGAT $-$ Naive & $-$0.040 & [$-$0.108, 0.023] & 0.828 & 0.850 & 0.838 & 0.828\,/\,1.000 \\
 & H2a: MGAT $-$ Graph-free & $-$0.016 & [$-$0.061, 0.029] & 0.707 & 0.710 & 0.711 & 1.000 \\
 & H2b: MGAT $-$ Ridge & $-$0.011 & [$-$0.065, 0.041] & 0.627 & 0.622 & 0.636 & 1.000 \\
\midrule
Healthcare & H1: MGAT $-$ Naive & 0.043 & [$-$0.066, 0.164] & 0.275 & 0.267 & 0.262 & 0.549\,/\,0.280 \\
 & H2a: MGAT $-$ Graph-free & 0.048 & [$-$0.001, 0.112] & 0.093 & 0.116 & 0.094 & 0.280 \\
 & H2b: MGAT $-$ Ridge & 0.042 & [$-$0.013, 0.104] & 0.124 & 0.134 & 0.090 & 0.280 \\
\bottomrule\end{tabular}}
\end{table*}

\subsection{Results by agent and policy behaviour}\label{subsec:sector}

{\color{blue}Table~\ref{tab:byagent} reports the comparison between MGAT and naive agents for each algorithm, and Table~\ref{tab:policy} describes how the policies allocate. These comparisons are descriptive and are not corrected for multiplicity.

The per-agent differences are mixed in sign. The largest is that of SAC in healthcare (0.103, 90\% interval $[0.040, 0.170]$); TD3 differs by 0.140 in healthcare and $-0.180$ in banking, with intervals that span zero. The policy diagnostics show why the per-agent results should not be read as evidence that the agents exploit the forecasts. Across agents, the mean daily rank correlation between the weights of the MGAT agents and the MGAT forecast is $-0.041$ in banking and $-0.027$ in healthcare, and the weights of the MGAT and naive agents have a mean rank correlation with the realised tradable return of at most 0.015 in absolute value for every agent. The policies differ mainly in how far they move from equal weights: SAC and PPO stay close to equal weighting (active share 0.02--0.09), A2C moves further (0.11--0.15), and TD3 and DDPG deviate most (0.32--0.40). In healthcare, the MGAT SAC policies have an active share of 0.03, compared with 0.09 for the naive SAC policies, so their advantage reflects allocations closer to the equal-weight portfolio, which performed well in 2025, rather than a tilt towards the stocks that MGAT ranks highly.

The agents therefore learned allocations that are nearly insensitive to the forecast rows of the state. Two features of the setting make this plausible. The forecast signal is weak relative to daily return noise (IC$_2$ of 0.05 at best), and the reward of Eq.~\eqref{eq:rl_reward} rewards a tilt only through its effect on the next day's portfolio return, so 484 training days provide a noisy learning signal for a cross-sectional tilt. In addition, the development-year forecasts that enter the training states rank the tradable return with different strength in the two training years (Table~\ref{tab:forecast_years}), so the association between the forecast rows and subsequent returns is itself unstable during training.\par}

\begin{table*}[!t]\centering\small\color{blue}
\caption{Sharpe ratio in 2025 (F1) by agent: mean over five seeds for the naive and MGAT sources, and the paired difference with its 90\% block-bootstrap interval (block length 10).}
\label{tab:byagent}
\setlength{\tabcolsep}{4pt}
\resizebox{\ifdim\width>\textwidth\textwidth\else\width\fi}{!}{%
\begin{tabular}{lrrrrrrrr}\toprule
 & \multicolumn{4}{c}{Banking} & \multicolumn{4}{c}{Healthcare} \\ \cmidrule(lr){2-5}\cmidrule(lr){6-9}
Agent & Naive & MGAT & $\Delta$SR & 90\% interval & Naive & MGAT & $\Delta$SR & 90\% interval \\ \midrule
PPO & 0.52 & 0.50 & $-$0.012 & [$-$0.055, 0.030] & 1.17 & 1.16 & $-$0.015 & [$-$0.069, 0.033] \\
A2C & 0.55 & 0.54 & $-$0.008 & [$-$0.067, 0.043] & 1.25 & 1.24 & $-$0.010 & [$-$0.140, 0.128] \\
SAC & 0.60 & 0.62 & 0.021 & [$-$0.007, 0.047] & 1.12 & 1.22 & 0.103 & [0.040, 0.170] \\
TD3 & 0.70 & 0.52 & $-$0.180 & [$-$0.437, 0.051] & 1.08 & 1.22 & 0.140 & [$-$0.186, 0.501] \\
DDPG & 0.55 & 0.53 & $-$0.019 & [$-$0.179, 0.154] & 1.15 & 1.14 & $-$0.003 & [$-$0.212, 0.249] \\
\bottomrule\end{tabular}}
\end{table*}
\begin{table*}[!t]\centering\small\color{blue}
\caption{Behaviour of the 2025 policies (F1; mean over five seeds; descriptive, not part of the registered tests). Active share: half the $L_1$ distance between the daily target weights and equal weights (0 for equal weighting). Loading: mean daily Spearman correlation between the weights executed at the close of $t+1$ and the MGAT forecast $\hat q_t$; for the naive agents, which do not observe $\hat q_t$, it serves as a reference. Weight IC: mean daily Spearman correlation between the executed weights and the realised return $q_t$.}
\label{tab:policy}
\setlength{\tabcolsep}{3.5pt}
\resizebox{\ifdim\width>\textwidth\textwidth\else\width\fi}{!}{%
\begin{tabular}{lrrrrrrrrrrrr}\toprule
 & \multicolumn{6}{c}{Banking} & \multicolumn{6}{c}{Healthcare} \\ \cmidrule(lr){2-7}\cmidrule(lr){8-13}
 & \multicolumn{2}{c}{Active share} & \multicolumn{2}{c}{Loading} & \multicolumn{2}{c}{Weight IC} & \multicolumn{2}{c}{Active share} & \multicolumn{2}{c}{Loading} & \multicolumn{2}{c}{Weight IC} \\
Agent & Naive & MGAT & Naive & MGAT & Naive & MGAT & Naive & MGAT & Naive & MGAT & Naive & MGAT \\ \midrule
PPO & 0.08 & 0.06 & $-$0.010 & 0.005 & 0.007 & $-$0.001 & 0.07 & 0.07 & $-$0.004 & 0.000 & $-$0.004 & $-$0.002 \\
A2C & 0.11 & 0.12 & 0.043 & $-$0.025 & 0.001 & 0.002 & 0.12 & 0.15 & 0.037 & $-$0.064 & 0.002 & $-$0.003 \\
SAC & 0.02 & 0.04 & $-$0.010 & $-$0.036 & 0.015 & 0.011 & 0.09 & 0.03 & 0.041 & $-$0.014 & $-$0.004 & 0.005 \\
TD3 & 0.37 & 0.37 & $-$0.001 & $-$0.080 & 0.008 & $-$0.003 & 0.40 & 0.32 & 0.093 & $-$0.055 & 0.003 & 0.009 \\
DDPG & 0.34 & 0.34 & $-$0.029 & $-$0.071 & 0.001 & 0.001 & 0.34 & 0.38 & 0.006 & $-$0.005 & 0.003 & $-$0.005 \\
\midrule
All & 0.18 & 0.19 & $-$0.001 & $-$0.041 & 0.007 & 0.002 & 0.20 & 0.19 & 0.035 & $-$0.027 & 0.000 & 0.001 \\
\bottomrule\end{tabular}}
\end{table*}

\subsection{Transaction-cost sensitivity}\label{subsec:cost}

{\color{blue}Table~\ref{tab:cost} reports the comparison between the MGAT and naive agents under all seven cost settings. Because the policies trade 4.6--9.7\% of the portfolio per day, higher costs lower all Sharpe ratios: from F0 to F4 the mean Sharpe ratio of the MGAT agents decreases from 0.61 to 0.38 in banking and from 1.25 to 1.08 in healthcare. Across the eight combinations of forecast source and sector, moving from F0 to the baseline setting F1 lowers the mean annualised return by a median of 1.8 percentage points (range 0.9--2.5), lowers the mean Sharpe ratio by 0.06 (0.05--0.09), and raises the mean maximum drawdown by 0.2 percentage points; at a 20~bp commission (F4), these changes are 5.8 (3.0--7.9) percentage points, 0.20 (0.17--0.28), and 0.6 (0.4--0.7) percentage points. The Sharpe ratio of each of the 200 policies decreases monotonically from F0 to F4. The differences between sources change little with costs (banking $-0.041$ to $-0.037$; healthcare 0.033 to 0.064). The healthcare difference increases slightly with costs because the MGAT agents trade less (8.4\% per day, compared with 9.7\% for the naive agents), but every interval includes zero.\par}

\begin{table}[!t]\centering\small\color{blue}
\caption{Cost sensitivity in 2025: mean Sharpe ratio of the 25 MGAT and naive policies and their paired difference (90\% block-bootstrap interval) under each cost setting. Commission/slippage in bp: F0 0/0 without stamp duty, F1 3/2, F2 5/2, F3 10/2, F4 20/2, F5 3/0, F6 3/5.}
\label{tab:cost}
\setlength{\tabcolsep}{3pt}
\begin{tabular}{lrrrrrr}\toprule
 & \multicolumn{3}{c}{Banking} & \multicolumn{3}{c}{Healthcare} \\ \cmidrule(lr){2-4}\cmidrule(lr){5-7}
Cost & MGAT & Naive & $\Delta$SR [90\%] & MGAT & Naive & $\Delta$SR [90\%] \\ \midrule
F0 & 0.61 & 0.65 & $-$0.041 [$-$0.110, 0.022] & 1.25 & 1.21 & 0.033 [$-$0.074, 0.154] \\
F1 & 0.54 & 0.58 & $-$0.040 [$-$0.108, 0.023] & 1.20 & 1.15 & 0.043 [$-$0.066, 0.164] \\
F2 & 0.52 & 0.56 & $-$0.039 [$-$0.108, 0.024] & 1.18 & 1.14 & 0.045 [$-$0.063, 0.167] \\
F3 & 0.48 & 0.52 & $-$0.038 [$-$0.107, 0.024] & 1.15 & 1.09 & 0.051 [$-$0.058, 0.174] \\
F4 & 0.38 & 0.42 & $-$0.037 [$-$0.105, 0.026] & 1.08 & 1.01 & 0.064 [$-$0.047, 0.188] \\
F5 & 0.56 & 0.60 & $-$0.040 [$-$0.108, 0.023] & 1.21 & 1.17 & 0.040 [$-$0.068, 0.161] \\
F6 & 0.51 & 0.55 & $-$0.039 [$-$0.108, 0.024] & 1.17 & 1.13 & 0.046 [$-$0.062, 0.168] \\
\bottomrule\end{tabular}
\end{table}

\subsection{Forecast-source ablation}\label{subsec:ablation}

{\color{blue}Table~\ref{tab:ablation_ppo} compares PPO agents (seeds 1--5) that observe forecasts from the alternative graph encoders. The Sharpe ratios of the seven sources lie within 0.50--0.55 in banking and 1.15--1.23 in healthcare. MGAT exceeds none of them by more than 0.005, and the GCN, graph-free, and ridge sources in banking and the no-cross-stock source in healthcare have higher Sharpe ratios, with 90\% intervals of the paired differences that exclude zero. With PPO, which stays close to equal weighting (Table~\ref{tab:policy}), the choice of graph encoder therefore has little effect on portfolio outcomes, in line with the small differences in IC$_2$ between the encoders.\par}

\begin{table}[!t]\centering\small\color{blue}
\caption{Forecast-source ablation with PPO (seeds 1--5), 2025, F1: mean Sharpe ratio of each source and the paired difference MGAT minus source (90\% block-bootstrap interval).}
\label{tab:ablation_ppo}
\setlength{\tabcolsep}{3pt}
\begin{tabular}{lrrrr}\toprule
 & \multicolumn{2}{c}{Banking} & \multicolumn{2}{c}{Healthcare} \\ \cmidrule(lr){2-3}\cmidrule(lr){4-5}
Source & SR & MGAT $-$ source & SR & MGAT $-$ source \\ \midrule
MGAT & 0.50 & -- & 1.16 & -- \\
GCN & 0.55 & $-$0.041 [$-$0.071, $-$0.013] & 1.18 & $-$0.022 [$-$0.078, 0.025] \\
No cross-stock & 0.51 & $-$0.005 [$-$0.039, 0.028] & 1.23 & $-$0.069 [$-$0.116, $-$0.023] \\
MGNN & 0.50 & 0.002 [$-$0.038, 0.042] & 1.16 & $-$0.001 [$-$0.072, 0.067] \\
Graph-free & 0.53 & $-$0.030 [$-$0.059, $-$0.003] & 1.20 & $-$0.044 [$-$0.106, 0.008] \\
Ridge & 0.54 & $-$0.035 [$-$0.065, $-$0.007] & 1.15 & 0.005 [$-$0.048, 0.055] \\
Naive & 0.52 & $-$0.012 [$-$0.055, 0.030] & 1.17 & $-$0.015 [$-$0.069, 0.033] \\
\bottomrule\end{tabular}
\end{table}

\subsection{Using the forecasts directly}\label{subsec:rules}

{\color{blue}To separate the information in the forecasts from the policy learner, Table~\ref{tab:forecast_rules} reports rules that hold, in equal weights, the three stocks with the highest $\hat q_t$, executed at the close of $t+1$. These rules are concentrated and replace 66--142\% of the account per day.

Without costs, the MGAT rule in healthcare attains a Sharpe ratio of 1.48 and a maximum drawdown of 12.8\%, compared with 1.25 and 19.1\% for daily equal weighting; the ridge and no-cross-stock rules attain 1.47 and 1.64. In banking, the MGAT rule attains 0.72, compared with 0.62. Under F1, the turnover costs reverse these differences: the MGAT rule attains 1.04 in healthcare and $-0.54$ in banking. In an exploratory block bootstrap, the frictionless difference in healthcare between the MGAT rule and equal weighting (0.23) has a 90\% interval of $[-0.73, 1.04]$, which reflects the variability of a concentrated three-stock portfolio over one year. The forecasts therefore carry economic information that a ranking rule can extract before costs, but extracting it requires more turnover than proportional costs of a few basis points allow.\par}

\begin{table*}[!t]\centering\small\color{blue}
\caption{Forecast-ranking rules in 2025: equal weights on the three stocks with the highest predicted return from the close of $t+1$ to the close of $t+2$, executed at the close of $t+1$, compared with daily equal weighting.}
\label{tab:forecast_rules}
\setlength{\tabcolsep}{4pt}
\resizebox{\ifdim\width>\textwidth\textwidth\else\width\fi}{!}{%
\begin{tabular}{llrrrrrr}\toprule
 & & \multicolumn{3}{c}{F0 (no costs)} & \multicolumn{3}{c}{F1} \\ \cmidrule(lr){3-5}\cmidrule(lr){6-8}
Sector & Strategy & AR (\%) & SR & MDD (\%) & AR (\%) & SR & MDD (\%) \\ \midrule
Banking & Equal weight, daily & 8.6 & 0.62 & 15.5 & 8.4 & 0.61 & 15.5 \\
 & MGAT & 11.4 & 0.72 & 18.2 & $-$10.1 & $-$0.54 & 24.7 \\
 & GCN & 19.9 & 1.27 & 13.8 & $-$0.5 & 0.05 & 17.3 \\
 & No cross-stock & 15.7 & 0.97 & 19.3 & $-$7.9 & $-$0.42 & 26.7 \\
 & MGNN & 12.9 & 0.83 & 17.7 & $-$3.9 & $-$0.17 & 20.3 \\
 & Graph-free & 1.3 & 0.16 & 23.3 & $-$13.5 & $-$0.71 & 26.2 \\
 & Ridge & 6.7 & 0.47 & 19.0 & $-$16.2 & $-$0.97 & 23.5 \\
\midrule
Healthcare & Equal weight, daily & 38.1 & 1.25 & 19.1 & 37.6 & 1.23 & 19.1 \\
 & MGAT & 45.9 & 1.48 & 12.8 & 28.7 & 1.04 & 15.0 \\
 & GCN & 36.9 & 1.22 & 19.2 & 16.7 & 0.68 & 21.8 \\
 & No cross-stock & 59.4 & 1.64 & 18.4 & 39.0 & 1.21 & 20.0 \\
 & MGNN & 26.2 & 0.94 & 15.9 & 5.3 & 0.32 & 21.2 \\
 & Graph-free & 33.6 & 1.19 & 17.7 & 17.6 & 0.73 & 18.8 \\
 & Ridge & 49.4 & 1.47 & 14.9 & 14.2 & 0.59 & 18.3 \\
\bottomrule\end{tabular}}
\end{table*}

\subsection{Post-hoc analyses requested in review}\label{subsec:posthoc}

{\color{blue}The analyses in this subsection were requested by the reviewers and specified after the registered evaluation, in two dated addenda to the protocol that were committed to the code repository before they were run. They are descriptive, are not part of the registered tests, and are reported in full; none of them changes the registered forecaster, agents, selection, or tests.

\paragraph{Relation layers and unrestricted inter-stock attention}
Table~\ref{tab:xablation} reports MGAT variants that remove one relation layer from the edge set of Eq.~\eqref{eq:Afused_re} or admit all ordered stock pairs, trained in the same four folds with three seeds. Because the positive entries of the Kendall graph alone cover 90\% (banking) and 82\% (healthcare) of the ordered pairs, and every shareholder edge is also a Kendall edge, removing the shareholder or Q\&A layer changes few or no admissible pairs: the variant without the shareholder layer has exactly the neighbourhoods of MGAT in both sectors, and so does the variant without the Q\&A layer in healthcare. These retrainings use the same code, data, and seeds as the registered runs but a different number of CPU threads, and therefore a different order of floating-point operations; they measure the run-to-run variation of the training itself: their IC$_2$ differs from that of the registered MGAT by up to 0.048 in a development year (banking, 2023) and by 0.006 or less in 2025. Removing the Kendall layer restricts attention to firms linked by shareholding or co-mentions (75\% and 36\% of the pairs); this variant has the highest 2025 IC$_2$ of the variants in both sectors (0.021 and 0.064) but not in the development years, and unrestricted attention over all pairs has a lower IC$_2$ than MGAT in banking and a slightly higher one in healthcare. In 2025, all paired daily IC$_2$ differences between MGAT and these variants lie between $-0.012$ and 0.011, with 95\% block-bootstrap intervals that include zero, and the high and low forecasts differ by less than one percentage point. The relation layers therefore mainly exclude a minority of stock pairs, and neither removing a layer nor admitting all pairs changes the forecasts beyond the variation between training runs. Unrestricted attention has the same number of parameters as MGAT and adds only the coefficients of the excluded pairs; its mean training time per run in fold 2025 (61--77~s, under a heavier machine load than the registered runs) is comparable with that of MGAT (62--65~s), so its computational overhead is negligible for fifteen stocks.

\paragraph{Graph-attention coefficients}
Table~\ref{tab:attention} summarises the attention coefficients of the registered MGAT on the 2025 inputs. In banking, the coefficients are close to uniform over the admissible neighbours in all three layers (entropy 94--100\% of that of uniform attention), and pairs linked by shareholding or co-mentions receive the same attention as other pairs (ratios 0.97--1.00). In healthcare, the first layer gives pairs co-mentioned in investor questions 24\% and pairs with common shareholders 10\% more weight than the average admissible neighbour and pairs linked only by trading-value correlation 6\% less, whereas the second and third layers are nearly uniform. The first-layer coefficients vary from day to day (median coefficient of variation 0.25 and 0.33) and agree only moderately across seeds (correlation 0.52 and 0.41), whereas those of the deeper layers are stable. The learned attention therefore reflects the qualitative relations weakly, and only in the first layer and in the sector where the forecasts carry ranking information; it does not single out a small set of economically interpretable links.

\paragraph{Depth of the graph and Transformer components}
Table~\ref{tab:xablation} also reports variants with one graph-attention layer (64 units; 42{,}526 parameters), two layers (32 and 64 units; 44{,}414 parameters), and two Transformer encoder blocks (84{,}790 parameters). In 2025, their IC$_2$ is 0.002--0.015 in banking and 0.044--0.065 in healthcare, against 0.014 and 0.052 for MGAT; the paired daily differences from MGAT lie between $-0.013$ and 0.012, with 95\% intervals that include zero, and their high and low forecasts differ from those of MGAT by at most 2.3 percentage points. Over the development years, the three-layer MGAT has the highest IC$_2$ in banking (0.057, against 0.002--0.041), and the two-layer variant the highest in healthcare (0.026, against 0.018 for MGAT). The depth of the graph encoder therefore does not change the conclusions, and a second Transformer block, which doubles the number of parameters, does not improve the forecasts.

\paragraph{Direction of the forecasts}
As a further decision-relevant measure, we computed the share of stock-days on which the sign of $\hat q_{t,i}$, or of $\hat q_{t,i}$ relative to the day's cross-sectional mean, matches that of the realised $q_{t,i}$. In 2025, MGAT has hit rates of 50.8\% and 50.2\% in banking and 50.5\% and 51.4\% in healthcare. Across all forecasters of Tables~\ref{tab:forecast_eval} and~\ref{tab:xablation}, the hit rates lie between 48.4\% and 52.6\% and the relative hit rates between 48.9\% and 52.1\%; the largest deviation from 50\% has $t=2.0$ (one-layer variant, healthcare, relative sign), and the other 51 hit rates have $|t|\le1.9$. Because the cross-sectional rank correlation uses the full ordering rather than signs, it is the more informative measure of the weak signal.

\paragraph{Allocation baselines}
Panel~A of Table~\ref{tab:posthoc_port} adds long-only allocation rules that use only the 252-day covariance matrix of the RL state, with the same ledger and costs. In banking, the minimum-variance portfolio attains a Sharpe ratio of 1.27 (monthly rebalancing) and 1.25 (daily) in 2025 under F1, with an annualised return of 19.9\% and 19.3\%, compared with 0.61 for daily equal weighting and 0.54--0.58 for the RL agents of every forecast source; it also exceeded daily equal weighting and the mean of the selected policies in the selection year 2024 (1.85, against 1.70 and 1.64--1.69). In healthcare, minimum variance (1.19) and inverse volatility (1.20) are close to equal weighting (1.23) and to the MGAT agents (1.20). A rule that requires no forecast thus outperformed every learned policy in banking in 2025. Because the agents observe the same covariance matrix, this result shows that they also did not exploit the risk information in their state; their allocations remain close to equal weighting (Table~\ref{tab:policy}).

\paragraph{Volatility regimes within 2025}
Panel~B splits the 2025 trading days at the median of the trailing 20-day volatility of daily equal weighting (121 days each). The paired difference between the MGAT and naive agents is $-0.008$ $[-0.115, 0.087]$ on high-volatility and $-0.068$ $[-0.144, 0.004]$ on low-volatility days in banking, and $-0.011$ $[-0.135, 0.133]$ and 0.107 $[-0.065, 0.267]$ in healthcare. No interval excludes zero; the positive healthcare difference of the full year comes from the calmer half of the year.

\paragraph{PPO learning rate and discount factor}
Panel~C of Table~\ref{tab:posthoc_port} retrains the PPO agents of the MGAT and naive sources (seeds 1--5) with one hyperparameter changed and otherwise the registered training and selection. Across the five settings, the mean Sharpe ratio of the MGAT agents lies between 0.48 and 0.57 in banking and between 1.16 and 1.21 in healthcare, against 0.46--0.57 and 1.14--1.20 for the naive agents. Seven of the eight changed settings give paired differences whose 90\% intervals include zero. With $\gamma=0.999$ in banking, the MGAT agents have a higher mean Sharpe ratio (0.066, $[0.020, 0.109]$, one-sided $p=0.008$), but the MGAT agent is higher for only two of the five seeds, and the difference is driven by large gaps in a minority of pairs (the standard deviation of the naive agents' Sharpe ratios is 0.13, against 0.03 for the MGAT agents); it is not significant after a Bonferroni correction for the eight comparisons (threshold 0.006). The conclusions of the registered PPO comparison therefore do not depend on the learning rate or the discount factor within these ranges.

\paragraph{Market structure}
Daily returns co-move strongly in banking: the mean pairwise correlation is 0.56--0.63 in 2022--2025, and the first principal component explains 60--65\% of the standardised variance. In healthcare, the correlation is 0.30--0.51, the first component explains 35--55\%, and the mean cross-sectional standard deviation of daily returns is 1.8--2.4\%, compared with 0.7--0.8\% in banking. The larger cross-sectional dispersion leaves more room for ranking stocks, which is consistent with the significant 2025 IC$_2$ in healthcare and its absence in banking; with two sectors, this association is descriptive.\par}

\begin{table*}[!t]\centering\small\color{blue}
\caption{Post-hoc forecaster ablation (not part of the registered tests). Each variant differs from MGAT only in the stated respect and is trained in the same four folds with seeds 1--3. IC$_2$: mean daily rank correlation of the predicted and realised tradable return (Eq.~\eqref{eq:ic2}), with its $t$-statistic for 2025; 2022--23: mean of the two development years. High, Low: change in the RMSE of the next-day log high and low in 2025 relative to the previous-day values (\%). Pairs: admissible ordered stock pairs (\%).}
\label{tab:xablation}
\setlength{\tabcolsep}{3pt}
\resizebox{\ifdim\width>\textwidth\textwidth\else\width\fi}{!}{%
\begin{tabular}{lrrrrrrrrrrrr}\toprule
 & \multicolumn{6}{c}{Banking} & \multicolumn{6}{c}{Healthcare} \\ \cmidrule(lr){2-7}\cmidrule(lr){8-13}
Forecaster & Pairs & 2022--23 & 2024 & 2025 ($t$) & High & Low & Pairs & 2022--23 & 2024 & 2025 ($t$) & High & Low \\ \midrule
MGAT (registered runs) & 92 & 0.057 & $-$0.007 & 0.014 (0.7) & $-$18.9 & $-$18.9 & 82 & 0.018 & 0.005 & 0.052 (2.3) & $-$15.5 & $-$28.5 \\
without Kendall layer & 75 & 0.048 & 0.008 & 0.021 (1.1) & $-$18.9 & $-$18.9 & 36 & 0.011 & 0.015 & 0.064 (3.2) & $-$15.3 & $-$28.2 \\
without Q\&A layer & 90 & 0.029 & 0.004 & 0.012 (0.6) & $-$18.0 & $-$18.5 & 82 & 0.020 & 0.011 & 0.051 (2.3) & $-$15.4 & $-$28.6 \\
without shareholder layer & 92 & 0.032 & 0.002 & 0.008 (0.4) & $-$18.1 & $-$18.5 & 82 & 0.020 & 0.011 & 0.051 (2.3) & $-$15.4 & $-$28.6 \\
all pairs admissible & 100 & 0.016 & 0.000 & 0.003 (0.1) & $-$18.2 & $-$18.9 & 100 & 0.017 & 0.051 & 0.057 (2.5) & $-$15.4 & $-$28.0 \\
one GAT layer (64) & 92 & 0.041 & 0.013 & 0.002 (0.1) & $-$19.7 & $-$17.1 & 82 & 0.012 & 0.035 & 0.065 (3.0) & $-$15.5 & $-$29.3 \\
two GAT layers (32, 64) & 92 & 0.002 & 0.012 & 0.015 (0.7) & $-$16.6 & $-$18.0 & 82 & 0.026 & 0.016 & 0.064 (3.1) & $-$15.0 & $-$29.2 \\
two Transformer blocks & 92 & 0.034 & 0.036 & 0.008 (0.4) & $-$19.7 & $-$20.3 & 82 & 0.013 & 0.019 & 0.044 (2.0) & $-$15.0 & $-$28.9 \\
\bottomrule\end{tabular}}
\end{table*}
\begin{table}[!t]\centering\small\color{blue}
\caption{Post-hoc analysis of the graph-attention coefficients of MGAT (fold 2025, three seeds, last day of every 2025 input window). Self: mean weight on the stock itself (uniform attention in parentheses). SH, QA, KD only: mean attention on admissible pairs linked by the shareholder graph, by the Q\&A graph, or only by the Kendall graph, relative to the mean over all admissible pairs (1 = no preference). Entropy: relative to uniform attention. CV: median day-to-day coefficient of variation of the coefficients. Seeds: mean correlation of the average coefficients between seeds.}
\label{tab:attention}
\setlength{\tabcolsep}{3pt}
\begin{tabular}{llrrrrrrr}\toprule
Sector & Layer & Self & SH & QA & KD only & Entropy & CV & Seeds \\ \midrule
Banking & 1 & 0.075 (0.074) & 1.00 & 0.98 & 1.01 & 0.94 & 0.25 & 0.52 \\
 & 2 & 0.074 (0.074) & 0.97 & 0.99 & 1.06 & 1.00 & 0.02 & 0.99 \\
 & 3 & 0.074 (0.074) & 0.97 & 0.99 & 1.06 & 1.00 & 0.00 & 1.00 \\
\midrule
Healthcare & 1 & 0.080 (0.087) & 1.10 & 1.24 & 0.94 & 0.82 & 0.33 & 0.41 \\
 & 2 & 0.089 (0.087) & 0.95 & 0.94 & 1.04 & 1.00 & 0.03 & 0.98 \\
 & 3 & 0.087 (0.087) & 0.95 & 0.94 & 1.04 & 1.00 & 0.01 & 1.00 \\
\bottomrule\end{tabular}
\end{table}
\begin{table*}[!t]\centering\small\color{blue}
\caption{Post-hoc portfolio analyses in 2025 (not part of the registered tests). Panel A: allocation baselines with the same ledger and costs (F0 without costs; F1 baseline costs); minimum variance uses the 252-day covariance of the RL state. Panel B: Sharpe ratios on the days with above- and below-median trailing 20-day volatility of daily equal weighting, mean over the 25 MGAT and 25 naive policies (F1), and the paired difference with its 90\% block-bootstrap interval. Panel C: PPO agents (seeds 1--5, F1) with one hyperparameter changed; registered: learning rate $10^{-3}$, $\gamma=0.99$.}
\label{tab:posthoc_port}
\setlength{\tabcolsep}{3pt}
\resizebox{\ifdim\width>\textwidth\textwidth\else\width\fi}{!}{%
\begin{tabular}{lrrrrrrrrrr}\toprule
\multicolumn{11}{l}{\textit{Panel A: allocation baselines}} \\
 & \multicolumn{5}{c}{Banking} & \multicolumn{5}{c}{Healthcare} \\ \cmidrule(lr){2-6}\cmidrule(lr){7-11}
Strategy & SR (F0) & AR (\%) & SR & MDD (\%) & Turn. (\%) & SR (F0) & AR (\%) & SR & MDD (\%) & Turn. (\%) \\ \midrule
Equal weight, daily & 0.62 & 8.4 & 0.61 & 15.5 & 1.0 & 1.25 & 37.6 & 1.23 & 19.1 & 1.8 \\
Minimum variance, monthly & 1.29 & 19.9 & 1.27 & 13.8 & 1.7 & 1.20 & 36.1 & 1.19 & 18.6 & 1.9 \\
Minimum variance, daily & 1.31 & 19.3 & 1.25 & 13.4 & 4.5 & 1.22 & 34.5 & 1.19 & 17.0 & 4.1 \\
Inverse volatility, monthly & 0.67 & 9.3 & 0.66 & 15.1 & 0.6 & 1.21 & 36.3 & 1.20 & 19.5 & 0.8 \\
MGAT agents (mean) & 0.61 & 7.4 & 0.54 & 15.6 & 5.6 & 1.25 & 36.8 & 1.20 & 18.9 & 8.4 \\
Naive agents (mean) & 0.65 & 8.1 & 0.58 & 15.9 & 5.8 & 1.21 & 34.7 & 1.15 & 19.6 & 9.7 \\
\midrule
\multicolumn{11}{l}{\textit{Panel B: volatility regimes within 2025}} \\
 & \multicolumn{5}{c}{Banking} & \multicolumn{5}{c}{Healthcare} \\ \cmidrule(lr){2-6}\cmidrule(lr){7-11}
Regime & EW & MGAT & Naive & $\Delta$SR & 90\% interval & EW & MGAT & Naive & $\Delta$SR & 90\% interval \\ \midrule
High volatility & 0.85 & 0.80 & 0.80 & $-$0.008 & [$-$0.115, 0.087] & 0.88 & 0.83 & 0.84 & $-$0.011 & [$-$0.135, 0.133] \\
Low volatility & 0.39 & 0.31 & 0.38 & $-$0.068 & [$-$0.144, 0.004] & 1.81 & 1.76 & 1.65 & 0.107 & [$-$0.065, 0.267] \\
\midrule
\multicolumn{11}{l}{\textit{Panel C: PPO hyperparameter sensitivity}} \\
 & \multicolumn{5}{c}{Banking} & \multicolumn{5}{c}{Healthcare} \\ \cmidrule(lr){2-6}\cmidrule(lr){7-11}
Setting & MGAT & Naive & $\Delta$SR & 90\% interval & Higher & MGAT & Naive & $\Delta$SR & 90\% interval & Higher \\ \midrule
Registered & 0.50 & 0.52 & $-$0.012 & [$-$0.055, 0.030] & 2/5 & 1.16 & 1.17 & $-$0.015 & [$-$0.069, 0.033] & 1/5 \\
Learning rate $3\times10^{-4}$ & 0.57 & 0.57 & $-$0.004 & [$-$0.022, 0.015] & 3/5 & 1.21 & 1.20 & 0.014 & [$-$0.012, 0.044] & 4/5 \\
Learning rate $3\times10^{-3}$ & 0.48 & 0.48 & 0.004 & [$-$0.089, 0.095] & 3/5 & 1.20 & 1.16 & 0.039 & [$-$0.067, 0.151] & 5/5 \\
$\gamma=0.9$ & 0.51 & 0.50 & 0.006 & [$-$0.033, 0.044] & 3/5 & 1.17 & 1.16 & 0.011 & [$-$0.067, 0.089] & 1/5 \\
$\gamma=0.999$ & 0.53 & 0.46 & 0.066 & [0.020, 0.109] & 2/5 & 1.18 & 1.14 & 0.038 & [$-$0.019, 0.108] & 2/5 \\
\bottomrule\end{tabular}}
\end{table*}


\subsection{Development of the forecaster}\label{subsec:history}

{\color{blue}The forecaster evaluated here is the outcome of a documented sequence of designs, all of which are archived in the code repository. The audit made during this revision showed that the forecaster of the original submission used valuation ratios of the forecast day as inputs. After correction, the multi-relation encoder of the original submission (MGNN) forecasted high and low prices more accurately than previous-day values but had no rank information at the tradable horizon, and agents whose training observations were built from its forecasts did not differ significantly from agents trained on observed prices in nine of ten agent--sector comparisons. We then trained all encoders with the ranking-aware objective of Section~\ref{subsec:objective} and fixed, before training, a gate on the development years 2022--2023 that required the MGNN encoder to rank the tradable return significantly and at least as well as the graph-free and linear comparators. The MGNN encoder failed this gate in both sectors (IC$_2$ $-0.011$ in banking and 0.012 in healthcare), and the evaluation of that design was stopped without using later data. Among the encoders that propagate information across stocks, graph attention ranked the tradable return best in both sectors, and it became the forecaster of the present protocol (Section~\ref{subsec:objective}). The 2025 evaluation was registered after these development steps and before the 2025 data were collected.\par}

\subsection{Computational cost}\label{subsec:compute}

{\color{blue}All experiments ran on one laptop CPU (Intel Core i7-12700H, 14 cores, 20 threads) without a GPU. Training one MGAT forecaster takes 1--3 minutes; the 15 forecasters of one sector and fold (five encoders, three seeds) take 17--23 minutes with nine threads. With 16 single-threaded RL processes in parallel, the median training time per policy is about 1 minute for PPO and A2C, 25 minutes for TD3, 42 minutes for SAC, and 66 minutes for DDPG, dominated by the off-policy updates; the 230 policies required 94 CPU-hours. Evaluating one policy on one year takes less than one second, and a daily decision, including the forecast, takes a few milliseconds, so inference cost is negligible for daily rebalancing.\par}

\section{Limitations and practical applicability}\label{sec:limitations}

{\color{blue}\paragraph{Scope of the evidence}
The evaluation covers one year, two sectors of fifteen stocks each, and one market. The evaluation year was not used in any earlier analysis and was evaluated once under a registered protocol, but one year of daily returns limits the precision of Sharpe-ratio comparisons: the 90\% intervals of the registered tests have widths of 0.09--0.23, larger than the differences between forecast sources (at most 0.05). The universes were fixed in hindsight and all firms survived to the end of 2025, which introduces survivorship bias. The paired intervals reflect the resampling of trading days for the 25 agent--seed pairs of each source and condition on one forecaster ensemble per source (the mean of three forecaster seeds), so they do not include the variation that retraining the forecasters with other seeds would add. The forecaster and the policies are defined on a fixed set of stocks and cannot be applied to unseen stocks without retraining. The results therefore describe the behaviour of the framework in these sectors and this period rather than its expected performance elsewhere.

\paragraph{Data and accounting}
Prices are unadjusted and the ledger does not credit cash dividends, which lowers the measured returns of high-dividend stocks such as banks and affects all strategies within a sector alike. The relation graphs are static and use information up to mid-2023, so the forecasts that enter the RL training states (2022--2023) are retrospective with respect to the graphs, whereas the selection and evaluation years are not. The 2025 valuation ratios were rebuilt from a public source to match the definitions of the original panels and agree with them over 2024 (Section~\ref{sec5}); the original panels compute PS and PCF from the previous fiscal year's statements from the first trading day of each year, which anticipates annual reports by up to four months, and the 2025 rows follow the same convention so that the inputs are defined identically in training and evaluation. Execution at the observed close with proportional costs ignores market impact, board lots, and price limits, which matter most for the concentrated forecast-ranking rules.

\paragraph{Model design and market adaptation}
The relation graphs cover 82--92\% of the stock pairs, so they constrain graph attention only weakly; sparser or time-varying graphs, for example from rolling shareholding or news co-occurrence, may separate the encoders more clearly. Because every shareholder edge is also a Kendall edge, the shareholder graph adds no admissible pair in the edge-set formulation; using the relation weights as attention biases or edge features would let the qualitative relations enter the forecaster directly. The forecaster is held fixed during policy learning. Joint fine-tuning with portfolio rewards, cost-aware objectives, longer training histories, and state designs that present the forecast as a ranking rather than as additional rows of a covariance-based state may help the agents use a weak signal; these options were not evaluated and would require the same controlled protocol. Hyperparameters and budgets of the agents were fixed in advance and not tuned, and the agents were not given turnover penalties; the post-hoc analyses of Section~\ref{subsec:posthoc} vary the PPO learning rate and discount factor and the depth of the graph and Transformer components, but not the input window or the learning rate of the forecaster.

\paragraph{Practical use}
The forecaster is inexpensive to train and to run, and its high and low forecasts are consistently more accurate than previous-day values, which may be useful for execution and risk limits. Its tradable-horizon ranking was informative in healthcare in 2025 but varied across years and sectors. In this study, neither the DRL agents nor concentrated ranking rules turned the ranking into a net gain over equal weighting after costs, so practical use would require portfolio construction that limits turnover.\par}

\section{Conclusion}\label{sec6}

{\color{blue}This study proposed MGAT-FinRL, a framework that combines relation graphs built from shareholding, investor-relations, and trading-activity data with a graph-attention Transformer forecaster and DRL portfolio agents, and evaluated it under a protocol registered before the evaluation data were collected. Three findings emerge. First, the forecaster predicts next-day high and low prices 15--29\% more accurately than previous-day values in an unseen year, and its forecasts of the tradable-horizon return carry statistically significant ranking information in healthcare (IC$_2=0.052$, $p=0.012$), but not in banking. Second, the relation graphs add little to this information beyond what a Transformer trained with the same ranking-aware objective obtains without them. Third, DRL agents that observe the forecasts at every decision do not attain a significantly higher Sharpe ratio than agents without forecast information: the difference is 0.043 in healthcare and $-0.040$ in banking, and the learned allocations are nearly insensitive to the forecasts. A simple ranking rule extracts the forecast information before costs but not after proportional costs.

The evaluation therefore locates the constraint of the forecast-to-policy pipeline at the interface between a weak cross-sectional signal and the policy learner, rather than in the forecaster alone. Future work should test policy designs that are built to exploit such signals, including cost-aware objectives and joint training, on longer evaluation periods and broader universes, using the same pre-registered and paired evaluation.\par}

\section*{Declarations}
\subsection*{Funding}
This work was supported by the National Key R\&D Program (Nos. 2022YFA1004000 and 2022YFA1004001).

\subsection*{Competing interests}
The authors declare that they have no known competing financial interests or personal relationships that could have appeared to influence the work reported in this paper.

\subsection*{Ethics approval}
Not applicable.

\subsection*{Consent to participate}
Not applicable.

\subsection*{Consent for publication}
Not applicable.

\subsection*{Data and code availability}
{\color{blue}The source code for graph construction, forecasting, data download and ingestion, the backtesting ledger, RL training, and evaluation, together with the registered protocol and its dated amendments, configuration files, random seeds, and the scripts that generate every table and figure of this article, is available at \url{https://github.com/ren148056842/MGT-FinRL} on the branch \texttt{revision-b3}. The repository also contains the forecasts, training histories, and metadata of all forecasting runs, the 2024 validation and 2025 evaluation trajectories of all RL policies and reference strategies, the ingestion checks of the 2025 data, and the analysis outputs. The market data remain subject to the terms of their providers (Tushare, Eastmoney, and Wind) and can be regenerated with the included download scripts; the trained policy checkpoints are available from the corresponding author on request.\par}

\subsection*{Author contributions}
Xiaoping Ren contributed to conceptualization, software, data curation,
formal analysis, visualization, and preparation of the original draft.
Zhiping Chen contributed to methodology, supervision, resource
provision, project administration, and review and editing of the
manuscript. Both authors read and approved the final manuscript.

\subsection*{Use of AI-assisted technologies}
During manuscript preparation, the authors used an AI-assisted language
tool for language refinement and editorial support. All generated
content was critically reviewed and revised by the authors, who take
full responsibility for the final manuscript.
\FloatBarrier
\clearpage
\appendix
\section{Forecasting models: configurations and additional results}
\label{app:model_complexity}

{\color{blue}Table~\ref{tab:param_count} reports the number of trainable parameters of each forecaster, together with the mean number of training epochs, the mean epoch selected by early stopping, and the mean wall-clock training time over the six runs of fold 2025 (two sectors, three seeds), measured while two forecaster processes shared the laptop CPU described in Section~\ref{subsec:compute}. All neural forecasters use the same inputs, window, Transformer block, output layer, optimiser, and ranking-aware objective and differ only in the graph branch. Table~\ref{tab:forecast_years} reports forecast accuracy in each test year.\par}

\begin{table}[!h]
\centering\small
\color{blue}
\caption{Forecasters: graph branch, trainable parameters, mean number of epochs run, mean early-stopping epoch, and mean training time per run in fold 2025. The ridge regressions have $160$ coefficients and one intercept each.}
\label{tab:param_count}
\setlength{\tabcolsep}{4pt}
\begin{tabular}{llrrrr}
\toprule
Forecaster & Graph branch & Parameters & Epochs & Best epoch & Time (s) \\
\midrule
MGAT & graph attention on $\bar{\mathbf{A}}$ & 44{,}846 & 17.0 & 7.0 & 64 \\
GCN & fixed normalised $\bar{\mathbf{A}}$ & 44{,}622 & 24.8 & 14.8 & 87 \\
No cross-stock & same layers, $\mathbf{A}=\mathbf{I}$ & 44{,}622 & 23.2 & 13.2 & 77 \\
MGNN & $\mathbf{W}_M+\sum_r\alpha_r\mathbf{A}_r$ & 44{,}850 & 20.2 & 10.2 & 73 \\
Graph-free & none & 2{,}718 & 34.5 & 24.5 & 30 \\
Ridge (two targets) & none & $2\times161$ & -- & -- & $<$1 \\
\bottomrule
\end{tabular}
\end{table}

\begin{table*}[!t]\centering\small\color{blue}
\caption{Forecast accuracy by test year (walk-forward folds). Columns as in Table~\ref{tab:forecast_eval}; 2022--2023 are the development years, 2024 the selection year, and 2025 the evaluation year.}
\label{tab:forecast_years}
\setlength{\tabcolsep}{3.5pt}
\resizebox{\ifdim\width>\textwidth\textwidth\else\width\fi}{!}{%
\begin{tabular}{llrrrrrrrr}\toprule
 & & \multicolumn{4}{c}{Banking} & \multicolumn{4}{c}{Healthcare} \\ \cmidrule(lr){3-6}\cmidrule(lr){7-10}
Year & Forecaster & High & Low & Close & IC$_2$ ($t$) & High & Low & Close & IC$_2$ ($t$) \\ \midrule
2022 & MGAT & $-$7.0 & $-$12.5 & +2.0 & 0.033 (1.4) & $-$15.5 & $-$21.1 & +0.5 & 0.011 (0.6) \\
 & GCN & $-$9.6 & $-$15.4 & +1.1 & 0.025 (1.1) & $-$15.0 & $-$21.5 & +1.0 & 0.027 (1.4) \\
 & No cross-stock & $-$9.9 & $-$16.4 & +1.3 & 0.040 (1.9) & $-$15.0 & $-$21.3 & +0.8 & 0.012 (0.7) \\
 & MGNN & $-$5.8 & $-$15.9 & +2.2 & $-$0.045 ($-$2.0) & $-$14.8 & $-$21.5 & +1.3 & 0.021 (1.1) \\
 & Graph-free & $-$9.8 & $-$16.4 & +0.3 & 0.055 (2.6) & $-$15.3 & $-$21.1 & +0.0 & 0.018 (1.0) \\
 & Ridge & -- & -- & -- & 0.061 (2.3) & -- & -- & -- & 0.005 (0.3) \\
\midrule
2023 & MGAT & $-$13.8 & $-$24.3 & +1.7 & 0.081 (3.5) & $-$14.2 & $-$20.5 & +1.4 & 0.026 (1.4) \\
 & GCN & $-$13.2 & $-$25.2 & +1.7 & 0.082 (3.6) & $-$13.6 & $-$20.0 & +3.1 & $-$0.007 ($-$0.3) \\
 & No cross-stock & $-$14.2 & $-$24.0 & +3.5 & 0.073 (3.2) & $-$14.6 & $-$21.3 & +0.6 & $-$0.009 ($-$0.5) \\
 & MGNN & $-$15.4 & $-$24.7 & +3.2 & 0.022 (1.2) & $-$12.5 & $-$20.6 & +1.8 & 0.003 (0.2) \\
 & Graph-free & $-$15.1 & $-$23.1 & +0.4 & 0.084 (3.5) & $-$13.9 & $-$19.8 & $-$0.1 & 0.004 (0.2) \\
 & Ridge & -- & -- & -- & 0.036 (1.9) & -- & -- & -- & 0.039 (2.1) \\
\midrule
2024 & MGAT & $-$13.5 & $-$20.1 & +1.7 & $-$0.007 ($-$0.3) & $-$18.9 & $-$21.8 & +2.0 & 0.005 (0.2) \\
 & GCN & $-$13.5 & $-$19.6 & +1.2 & 0.015 (0.7) & $-$18.5 & $-$22.7 & +1.3 & 0.004 (0.2) \\
 & No cross-stock & $-$13.2 & $-$21.0 & +3.3 & 0.005 (0.2) & $-$19.3 & $-$23.2 & +0.1 & 0.026 (1.4) \\
 & MGNN & $-$14.9 & $-$20.1 & +1.1 & 0.000 (0.0) & $-$17.3 & $-$21.6 & +0.6 & 0.010 (0.5) \\
 & Graph-free & $-$13.4 & $-$19.7 & +0.3 & 0.031 (1.4) & $-$18.3 & $-$22.8 & +0.0 & 0.017 (0.8) \\
 & Ridge & -- & -- & -- & $-$0.030 ($-$1.2) & -- & -- & -- & 0.020 (1.0) \\
\midrule
2025 & MGAT & $-$18.9 & $-$18.9 & +6.1 & 0.014 (0.7) & $-$15.5 & $-$28.5 & +0.5 & 0.052 (2.3) \\
 & GCN & $-$19.7 & $-$18.1 & +5.4 & 0.025 (1.3) & $-$15.8 & $-$29.9 & +0.4 & 0.053 (2.4) \\
 & No cross-stock & $-$19.8 & $-$19.8 & +2.5 & 0.032 (1.7) & $-$15.7 & $-$28.2 & +1.2 & 0.044 (2.1) \\
 & MGNN & $-$14.9 & $-$18.9 & +4.1 & 0.000 (0.0) & $-$14.5 & $-$27.4 & +1.9 & 0.038 (1.8) \\
 & Graph-free & $-$20.2 & $-$20.4 & +0.2 & 0.027 (1.3) & $-$14.8 & $-$27.4 & +0.1 & 0.047 (2.1) \\
 & Ridge & -- & -- & -- & $-$0.030 ($-$1.6) & -- & -- & -- & 0.026 (1.2) \\
\bottomrule\end{tabular}}
\end{table*}
\FloatBarrier

\section{Training and validation convergence}
\label{app:convergence}
{\color{blue}Figure~\ref{fig:fig9} shows the training and validation losses (Eq.~\eqref{eq:loss}) of the five neural forecasters in fold 2025, for each seed. The validation loss of the graph-based models reaches its minimum after 7--15 epochs on average, after which the training loss continues to fall while the validation loss stays flat or rises; early stopping retains the weights at the minimum. The validation loss of the graph-free Transformer, which has about one-sixteenth of the parameters, decreases more gradually and reaches its minimum after 25 epochs on average.\par}

\begin{figure*}[!h]
\centering
\includegraphics[width=\textwidth]{fig9_b3.pdf}
\caption{{\color{blue}Training (dashed) and validation (solid) losses of the forecasters in fold 2025 (training targets up to 2023, validation year 2024); one line per seed; dots mark the early-stopping epoch.}}
\label{fig:fig9}
\end{figure*}
\FloatBarrier

{\color{blue}
\section{RL hyperparameters and training budgets}\label{app:hyper}
Table~\ref{tab:hyper} lists the settings used for all agents. They were fixed in advance, registered with the protocol, and implemented with Stable-Baselines3 2.4.0 and its default multilayer-perceptron policies; unlisted parameters take the library defaults. Checkpoints are saved at 20, 40, 60, 80, and 100\% of each budget and selected on 2024 (Section~\ref{subsec:splits}).

\begin{table}[!h]
\centering\small
\caption{RL hyperparameters and training budgets (environment steps over the 484-day training period). All agents use $\gamma=0.99$, reward scaling $\kappa_R=0.1$, an initial capital of $10^6$, and seeds 1--5.}
\label{tab:hyper}
\begin{tabular}{lp{8.2cm}r}
\toprule
Agent & Settings & Budget \\
\midrule
A2C & learning rate $4\times10^{-4}$, $n_{\mathrm{steps}}=10$, entropy coef.\ 0.005, GAE $\lambda=1.0$, RMSProp & 40{,}000 \\
PPO & learning rate $10^{-3}$, $n_{\mathrm{steps}}=2048$, batch 128, 10 epochs, clip 0.2, GAE $\lambda=0.95$, entropy coef.\ 0.005 & 40{,}960 \\
SAC & learning rate $10^{-4}$, batch 128, buffer $10^6$, learning starts 100, $\tau=0.005$, automatic entropy (initial 0.1) & 60{,}000 \\
TD3 & learning rate $10^{-3}$, batch 100, buffer $10^6$, $\tau=0.005$, policy delay 2, target noise 0.2 & 30{,}000 \\
DDPG & learning rate $10^{-3}$, batch 256, buffer $10^6$, $\tau=0.005$ & 50{,}000 \\
\bottomrule
\end{tabular}
\end{table}
\par}

{\color{blue}
\section{Stock universes}\label{app:universe}
\begin{table}[!h]
\centering\small
\caption{The two fixed stock universes.}
\label{tab:universe}
\setlength{\tabcolsep}{4pt}
\begin{tabular}{llll}
\toprule
\multicolumn{2}{l}{Banking} & \multicolumn{2}{l}{Healthcare innovation} \\
Code & Name & Code & Name \\
\midrule
000001.SZ & Ping An Bank (PAB) & 002044.SZ & Meinian Onehealth \\
600000.SH & SPDB & 002173.SZ & Chuangxin Medical \\
600015.SH & Hua Xia Bank (HXB) & 002219.SZ & Xinlicheng \\
600016.SH & China Minsheng Bank (CMBC) & 002524.SZ & Guangzheng Eye Hospital \\
600036.SH & China Merchants Bank (CMB) & 300015.SZ & Aier Eye Hospital \\
601166.SH & Industrial Bank (CIB) & 300143.SZ & Yingkang Life \\
601288.SH & Agricultural Bank of China (ABC) & 300149.SZ & Ruizhi Pharma \\
601328.SH & Bank of Communications (BOCOM) & 300244.SZ & Dian Diagnostics \\
601398.SH & ICBC & 300347.SZ & Tigermed \\
601658.SH & Postal Savings Bank (PSBC) & 300404.SZ & Boji Medical \\
601818.SH & China Everbright Bank (CEB) & 600645.SH & Zhongyuan Union \\
601916.SH & China Zheshang Bank (CZB) & 600721.SH & Baihua Pharma \\
601939.SH & China Construction Bank (CCB) & 600763.SH & Topchoice Medical \\
601988.SH & Bank of China (BOC) & 603127.SH & JOINN Laboratories \\
601998.SH & China CITIC Bank (ECITIC) & 603882.SH & KingMed Diagnostics \\
\bottomrule
\end{tabular}
\end{table}
\par}
\section{Portfolio performance metrics}\label{app:portfolio_metrics}

This appendix provides the detailed definitions of the portfolio performance metrics.
Let $\{V_t\}_{t=0}^{T}$ denote the account-value trajectory over the out-of-sample trading period, where $V_0$ is the initial capital and $V_T$ is the final account value.
The one-period portfolio return is defined as
\begin{equation}
r_t=\frac{V_t-V_{t-1}}{V_{t-1}}, \quad t=1,2,\ldots,T .
\end{equation}
In this paper, the annualization factor is set to $A=252$, corresponding to the number of trading days in a year.
The annual risk-free rate is denoted by $r_f$ and is set to zero in the empirical evaluation unless otherwise stated.

\subsection{Cumulative return}

Cumulative return (CR) measures the total return over the whole backtesting period:
\begin{equation}
\mathrm{CR}=\frac{V_T}{V_0}-1 .
\end{equation}
A larger CR indicates a higher total wealth increase during the test period.

\subsection{Annualised return}

Annualised return (AR) converts the cumulative return over $T$ trading days into an annualised return:
\begin{equation}
\mathrm{AR}=
\left(\frac{V_T}{V_0}\right)^{\frac{A}{T}}-1 .
\end{equation}
When the test period is exactly one trading year, AR is equal to CR.
Otherwise, AR provides a time-normalized measure of profitability.

\subsection{Annualised volatility}

The annualised volatility of portfolio returns is computed as
\begin{equation}
\sigma_{\mathrm{ann}}=\sqrt{A}\cdot \mathrm{Std}(r_t),
\end{equation}
where $\mathrm{Std}(r_t)$ denotes the standard deviation of daily portfolio returns.

\subsection{Sharpe ratio}

The Sharpe ratio (SR) measures the excess return per unit of total risk. \textcolor{blue}{It is computed from daily returns as
\begin{equation}
\mathrm{SR}=\sqrt{A}\,\frac{\overline{r_t-r_f^{d}}}{\mathrm{Std}(r_t-r_f^{d})},\qquad r_f^{d}=(1+r_f)^{1/A}-1,
\end{equation}
where the bar denotes the sample mean and Std uses the sample standard deviation.}
A higher Sharpe ratio indicates better risk-adjusted performance.

\subsection{Maximum drawdown}

Maximum drawdown (MDD) measures the largest peak-to-trough decline in the account value:
\begin{equation}
\mathrm{MDD}=
\max_{1\leq t \leq T}
\left(
\frac{\max_{0\leq s\leq t}V_s-V_t}
{\max_{0\leq s\leq t}V_s}
\right).
\end{equation}
\textcolor{blue}{MDD is reported as a non-negative fraction, so a smaller value indicates a smaller downside loss.}

\subsection{Calmar ratio}

The Calmar ratio measures annualised return relative to the magnitude of the maximum drawdown:
\begin{equation}
\mathrm{Calmar}=
\frac{\mathrm{AR}}{\mathrm{MDD}}.
\end{equation}
A higher Calmar ratio indicates a better trade-off between return generation and drawdown risk.
When AR is negative, the Calmar ratio can also be negative.

\subsection{Sortino ratio}

Different from the Sharpe ratio, the Sortino ratio only penalizes downside volatility.
The annualised downside deviation is defined as
\begin{equation}
\sigma_{\mathrm{down}}=
\sqrt{
A \cdot
\frac{1}{T}
\sum_{t=1}^{T}
\min \left(r_t-\frac{r_f}{A},0\right)^2
}.
\end{equation}
The Sortino ratio is then computed as
\begin{equation}
\mathrm{Sortino}=
\frac{\mathrm{AR}-r_f}{\sigma_{\mathrm{down}}}.
\end{equation}
a higher Sortino ratio indicates better downside-risk-adjusted performance.

\subsection{Final account value}

The final account value is the terminal portfolio wealth at the end of the backtesting period:
\begin{equation}
\mathrm{Final\ Value}=V_T .
\end{equation}
In our experiments, the initial capital is set to $V_0=1{,}000{,}000$.
Therefore, the final account value directly reflects the terminal wealth generated by each strategy.
\clearpage
\bibliographystyle{plainnat}
\bibliography{sn-bibliography-revised}

\end{document}
